\documentclass[12pt,a4paper]{article}

% ---------- Packages ----------
\usepackage[utf8]{inputenc}
\usepackage[T1]{fontenc}
\usepackage[french]{babel}
\usepackage{textcomp}
\usepackage{geometry}
\usepackage{enumitem}
\usepackage{xcolor}
\usepackage{titlesec}
\usepackage{array}
\usepackage{tabularx}
\usepackage{longtable}
\usepackage{booktabs}
\usepackage{tikz}
\usetikzlibrary{positioning}
\usepackage{hyperref}

\geometry{margin=2.5cm}
\setlength{\parindent}{0pt}
\setlength{\parskip}{6pt}
\renewcommand{\arraystretch}{1.25}

\titleformat{\section}
  {\Large\bfseries\color{blue!60!black}}
  {\thesection}{0.7em}{}

\titleformat{\subsection}
  {\large\bfseries}
  {\thesubsection}{0.7em}{}

\titleformat{\subsubsection}
  {\normalsize\bfseries\color{blue!40!black}}
  {\thesubsubsection}{0.7em}{}

\titleformat{\paragraph}[hang]
  {\normalsize\bfseries\color{blue!60!black}}
  {}{0pt}{}
\titlespacing*{\paragraph}{0pt}{10pt}{4pt}

% ---------- Commandes utiles ----------
% Bloc d'en-tête de chaque membre
\newcommand{\membre}[2]{%
  \textbf{Nom et prénom :} #1 \par
  \textbf{Questions :} #2 \par
}
% Rubriques du modèle de réponse
\newcommand{\reponse}{\textbf{Réponse :}\par}
\newcommand{\exemple}{\textbf{Exemple ou application pratique :}\par}
\newcommand{\sources}{\textbf{Source(s) :}}
% Flèche
\newcommand{\fl}{\ensuremath{\rightarrow}}
% Colonnes de tableaux alignées à gauche
\newcolumntype{L}[1]{>{\raggedright\arraybackslash}p{#1}}
\newcolumntype{Y}{>{\raggedright\arraybackslash}X}

% ---------- Informations du TP ----------
\newcommand{\matiere}{Services Réseaux}
\newcommand{\groupe}{Groupe de 14 étudiants}
\newcommand{\universite}{USTHB}
\newcommand{\annee}{2026--2027}

% ---------- Document ----------
\begin{document}

% ---------- Page de garde ----------
\begin{titlepage}
    \centering

    {\large République Algérienne Démocratique et Populaire \par}
    \vspace{0.2cm}
    {\large Ministère de l'Enseignement Supérieur et de la Recherche Scientifique \par}
    \vspace{0.6cm}
    {\Large \textbf{Université des Sciences et de la Technologie Houari Boumediene} (\universite) \par}
    \vspace{0.2cm}
    {\large Département d'Informatique \par}
    \vspace{1.5cm}

    {\Huge\bfseries TP 1 \par}
    \vspace{0.5cm}

    {\LARGE Découverte des Services Réseaux \par}
    \vspace{1.5cm}

    \rule{\textwidth}{0.5pt}

    \vspace{0.7cm}
    {\large Répartition du travail et réponses aux questions \par}
    \vspace{0.7cm}

    \rule{\textwidth}{0.5pt}

    \vspace{1.5cm}

    \textbf{Module :} ADV Réseau -- \matiere

    \textbf{Année universitaire :} \annee

    \vfill

    \textbf{Membres du groupe :} voir la répartition ci-dessous.

    \vfill
\end{titlepage}

\tableofcontents
\newpage

% ==================================================
\section{Consignes générales}
% ==================================================

Ce document présente la répartition du TP entre les
14 membres du groupe.

Chaque membre doit :
\begin{itemize}
    \item traiter les questions qui lui sont attribuées ;
    \item rédiger des réponses claires, structurées et précises ;
    \item ajouter une source fiable pour chaque réponse ;
    \item fournir des exemples ou des schémas lorsque nécessaire ;
    \item vérifier les informations avant de les intégrer au document final.
\end{itemize}

Les réponses doivent être rédigées avec les propres mots
des étudiants. Les sources utilisées doivent être citées
dans le texte et regroupées dans la bibliographie finale.

% ==================================================
\section{Répartition du travail et réponses}
% ==================================================

% --------------------------------------------------
\subsection{Membre 1 : DNS}
% --------------------------------------------------
\membre{\dotfill}{1 à 5}

\paragraph{A. DNS (Domain Name System)}

\subsubsection*{Question 1 (Facile) : Rôle principal du DNS}

\reponse
Le DNS traduit les noms de domaine lisibles par l'humain en adresses IP
exploitables par les machines. Il évite ainsi de mémoriser des adresses
numériques et permet de changer l'adresse d'un serveur sans modifier son
nom. Il fournit aussi d'autres informations via différents types
d'enregistrements.

\subsubsection*{Question 2 (Moyen) : Serveur DNS autoritaire, serveur DNS récursif et résolveur DNS}

\reponse
La différence entre un serveur DNS autoritaire, un serveur DNS récursif
et un résolveur DNS :

\begin{longtable}{|L{3.4cm}|L{11.6cm}|}
\hline
\textbf{Serveur DNS autoritaire} &
Le serveur DNS autoritaire est le serveur officiel qui détient les
véritables dossiers d'un nom de domaine spécifique. Détient les données
officielles d'une zone (ex. universite.dz) et répond avec autorité pour
les noms de cette zone. Il ne cherche pas ailleurs : le serveur
autoritaire répond uniquement sur ce qu'il gère (réponse définitive et
exacte). \\ \hline
\textbf{Serveur DNS récursif} &
Le serveur DNS récursif est un serveur distant (généralement géré par
notre fournisseur d'accès à Internet ou par des services publics).

Il reçoit la requête d'un client et fait tout le travail de résolution
à sa place en interrogeant et en parcourant toute la hiérarchie
d'Internet (les serveurs Racines, puis les serveurs TLD comme le .com ou
.fr) et à chaque recherche il conserve les résultats en cache pour ne
pas refaire tout le travail pendant une durée limitée. \\ \hline
\textbf{Résolveur DNS} &
Le résolveur DNS est un composant logiciel léger, généralement intégré
directement à notre système d'exploitation (Windows, macOS, Linux) ou à
notre navigateur web. Il sert d'intermédiaire de premier niveau.
Lorsqu'on tape une adresse web, il regarde d'abord dans son cache local.
S'il ne connaît pas la réponse, il formule la requête et l'envoie au
serveur DNS récursif configuré sur notre réseau (comme celui des box
internet). \\ \hline
\end{longtable}

\subsubsection*{Question 3 (Avancé) : Analyse du processus de résolution du nom www.universite.dz}

\reponse
\underline{Étape 1 : Vérification du cache local}

L'ordinateur (le résolveur local) vérifie d'abord son propre cache
système ou celui du navigateur web.

Résultat : si vous avez visité www.universite.dz récemment, l'adresse
IP est trouvée immédiatement sur votre machine. Zéro requête sur
Internet, la réponse est instantanée.

\underline{Étape 2 : L'interrogation du serveur Racine (Root)}

Action : n'ayant pas la réponse en mémoire, le serveur récursif initie
la recherche au sommet de la hiérarchie internet.

Serveurs impliqués : le serveur DNS récursif interroge l'un des 13
serveurs Racines (Root Servers, «~.~»).

Réponse : le serveur racine ne connaît pas l'IP finale, mais il renvoie
les adresses des serveurs gérant l'extension de premier niveau (TLD)
.dz.

\underline{Étape 3 : L'interrogation du serveur TLD .dz}

Action : le serveur récursif descend d'un niveau dans l'arbre DNS.

Serveurs impliqués : le serveur DNS récursif contacte l'un des serveurs
de premier niveau (TLD Server) chargés de la zone .dz (gérée en Algérie
par le CERIST).

Réponse : le serveur TLD .dz renvoie l'adresse du serveur DNS
autoritaire délégué officiellement pour gérer le domaine de second
niveau universite.dz.

\underline{Étape 4 : L'interrogation du serveur Autoritaire}

Action : le serveur récursif s'adresse directement à la source
officielle de l'information.

Serveurs impliqués : le serveur DNS récursif interroge le serveur DNS
autoritaire du domaine universite.dz (hébergé par l'université ou son
prestataire).

Réponse : ce serveur lit son fichier de zone, trouve l'enregistrement
de type A (IPv4) correspondant à l'hôte www, et renvoie l'adresse IPv4
finale.

\underline{Étape 5 : Transmission finale et mise en cache}

Action : la réponse effectue le chemin inverse pour clore le processus.

Serveurs impliqués : le serveur DNS récursif reçoit l'IP, en stocke une
copie dans sa mémoire cache (selon la durée définie par le TTL) pour
accélérer les futures demandes, puis transmet cette IP au résolveur
local du client.

Résultat : le navigateur du client utilise l'IP obtenue pour initier la
connexion vers le serveur web hébergeant le site de l'université.

\subsubsection*{Question 4 (Pratique) : Les enregistrements DNS qu'on utilise}

\reponse
\begin{itemize}
    \item Nom vers adresse IPv4 : \textbf{A}.
    \item Serveur de messagerie : \textbf{MX}.
    \item Alias vers un autre nom : \textbf{CNAME}.
    \item Nom vers adresse IPv6 : \textbf{AAAA}.
\end{itemize}

\subsubsection*{Question 5 (Réflexion) : Pourquoi DNSSEC est-il important ?}

\reponse
Le DNS a été conçu sans mécanisme d'authentification : un résolveur qui
reçoit une réponse n'a aucun moyen de vérifier qu'elle provient bien du
serveur légitime ni qu'elle n'a pas été modifiée en route. Un attaquant
peut donc fournir de fausses réponses, et l'utilisateur est redirigé
vers un site malveillant sans s'en apercevoir.

\begin{longtable}{|L{4.5cm}|L{10.5cm}|}
\hline
Empoisonnement de cache & Injection de fausses données dans le cache
d'un résolveur \\ \hline
Spoofing DNS & Usurpation de réponses DNS pour rediriger vers un faux
serveur \\ \hline
Attaque de l'homme du milieu (MITM) & Modification des réponses DNS en
transit \\ \hline
Redirection / pharming & Envoi des utilisateurs vers un faux site
(phishing, vol d'identifiants, malware) \\ \hline
\end{longtable}

\underline{Limites à connaître :}
\begin{itemize}
    \item DNSSEC ne chiffre pas les requêtes : il n'assure pas la
    confidentialité (pour cela : DoT ou DoH).
    \item Il ne protège pas contre les attaques DDoS et peut même
    augmenter la taille des réponses, ce qui favorise l'amplification.
    \item Il protège la résolution DNS, mais pas le site de destination
    lui-même (HTTPS/TLS reste nécessaire).
    \item Son déploiement est complexe (gestion et rotation des clés) et
    n'est pas encore universel.
\end{itemize}

\underline{Importance du DNSSEC :}

DNSSEC est important parce qu'il rend le DNS vérifiable, ce qui garantit
que l'on atteint bien la bonne destination. C'est un maillon essentiel
de la sécurité de l'infrastructure Internet.

\sources
\begin{itemize}
    \item IETF, RFC 1034 et RFC 1035 : Domain Names,
    \url{https://www.rfc-editor.org/rfc/rfc1035}.
    \item IETF, RFC 4033 : DNS Security Introduction and Requirements,
    \url{https://www.rfc-editor.org/rfc/rfc4033}.
\end{itemize}

% --------------------------------------------------
\subsection{Membre 2 : DHCP et NTP}
% --------------------------------------------------
\membre{\dotfill}{6 à 12}

\paragraph{B. DHCP (Dynamic Host Configuration Protocol)}

\subsubsection*{Question 6 (Facile) : Les 4 étapes du processus DHCP (DORA)}

\reponse
\begin{longtable}{|L{2.4cm}|L{2.9cm}|L{2.9cm}|L{6.4cm}|}
\hline
\textbf{Étape} & \textbf{Message} & \textbf{Sens} & \textbf{Explication} \\ \hline
\endhead
Discover & DHCPDISCOVER & Client \fl{} tous (broadcast) &
Le client, qui n'a pas encore d'adresse IP, diffuse une requête pour
trouver un serveur DHCP sur le réseau. \\ \hline
Offer & DHCPOFFER & Serveur \fl{} client &
Un serveur DHCP répond en proposant une adresse IP disponible, avec les
paramètres associés (masque, passerelle, DNS, durée du bail). \\ \hline
Request & DHCPREQUEST & Client \fl{} serveur (broadcast) &
Le client accepte l'offre reçue et demande officiellement cette adresse.
Le broadcast informe aussi les autres serveurs que leur offre est
refusée. \\ \hline
Acknowledge & DHCPACK & Serveur \fl{} client &
Le serveur confirme l'attribution : l'adresse est réservée au client
pour la durée du bail et le client peut l'utiliser. \\ \hline
\end{longtable}

\underline{Remarques :}
\begin{itemize}
    \item Les échanges utilisent UDP, port 67 côté serveur et port 68
    côté client.
    \item Au moment du Discover, le client utilise l'adresse source
    0.0.0.0 et l'adresse destination 255.255.255.255.
    \item Le bail (lease) a une durée limitée : le client le renouvelle
    généralement à 50~\% de sa durée (T1).
\end{itemize}

\subsubsection*{Question 7 (Moyen) : Informations fournies en plus de l'adresse IP}

\reponse
Un serveur DHCP peut distribuer, via des options DHCP, toute la
configuration réseau du client :
\begin{itemize}
    \item Masque de sous-réseau
    \item Passerelle par défaut (routeur)
    \item Serveurs DNS
    \item Nom de domaine / suffixe DNS
    \item Serveurs NTP
    \item Durée du bail (lease time)
    \item Serveur TFTP et nom du fichier de démarrage (boot PXE)
    \item Paramètres spécifiques à un fabricant (ex. : adresse du
    contrôleur Wi-Fi, du serveur de téléphonie IP)
    \item Liste de recherche de domaines et des serveurs WINS, des
    réservations (adresse fixe liée à une adresse MAC) et des paramètres
    de renouvellement (T1, T2).
\end{itemize}

\subsubsection*{Question 8 (Avancé) : Serveur DHCP sur un autre VLAN}

\reponse
\textbf{Problème :} le DHCPDISCOVER est un broadcast, qui reste confiné
au VLAN (domaine de diffusion) du client. Il ne peut donc pas atteindre
un serveur situé dans un autre VLAN.

\textbf{Solution :} le \textbf{DHCP Relay} (agent de relais), aussi
appelé \texttt{ip helper-address} sur Cisco.

\textbf{Fonctionnement :}
\begin{itemize}
    \item Le client diffuse son DHCPDISCOVER dans son VLAN.
    \item Le routeur ou switch L3 qui porte l'interface (SVI) de ce
    VLAN, configuré en relais DHCP, reçoit le broadcast.
    \item Il le convertit en unicast et le transmet directement au
    serveur DHCP, en renseignant le champ \textbf{giaddr} (Gateway IP
    Address) avec l'adresse IP de son interface côté client.
    \item Le serveur utilise giaddr pour savoir de quel sous-réseau
    provient la demande et choisit l'adresse dans le pool
    correspondant.
    \item La réponse (Offer, puis Ack) repart en unicast vers le relais,
    qui la retransmet au client.
\end{itemize}

\subsubsection*{Question 9 (Cas pratique) : Bail DHCP de 8 heures}

\reponse
À T1 = 50~\% du bail = 4 heures. Le client envoie un DHCPREQUEST en
unicast au serveur qui lui a attribué l'adresse. Si le serveur répond
par un DHCPACK, le bail est prolongé pour une nouvelle durée complète de
8 heures.

\textbf{Que se passe-t-il si le serveur ne répond pas ?}

\begin{longtable}{|L{3.2cm}|L{3.8cm}|L{7.6cm}|}
\hline
\textbf{Moment} & \textbf{Phase} & \textbf{Action du client} \\ \hline
4 h (T1) & Renouvellement (\textit{renewing}) & DHCPREQUEST en unicast
vers le serveur d'origine. Il réessaie régulièrement s'il n'obtient pas
de réponse. \\ \hline
7 h (T2 = 87,5~\%) & Rebinding & DHCPREQUEST en \textbf{broadcast} vers
n'importe quel serveur DHCP du réseau. \\ \hline
8 h & Expiration du bail & Le client \textbf{doit abandonner l'adresse}
et relancer le processus DORA depuis le début. \\ \hline
\end{longtable}

Si aucun serveur ne répond à l'expiration :
\begin{itemize}
    \item Le client perd sa configuration IP et ne peut plus communiquer
    normalement.
    \item Sous Windows, il s'attribue une adresse APIPA (169.254.x.x),
    utilisable uniquement en local.
\end{itemize}

\textbf{Conclusion :} T1 = 50~\% et T2 = 87,5~\% sont les valeurs par
défaut définies par la RFC 2131, mais le serveur peut les modifier via
les options 58 (T1) et 59 (T2).

\paragraph{C. NTP (Network Time Protocol)}

\subsubsection*{Question 10 (Facile) : Pourquoi la synchronisation temporelle est-elle critique dans un réseau d'entreprise ?}

\reponse
Le NTP (Network Time Protocol) est un protocole qui permet de
synchroniser les horloges des serveurs, des routeurs, des switches et
des autres équipements réseau. En entreprise, cette synchronisation est
essentielle pour garantir la cohérence des journaux (logs), faciliter le
dépannage, renforcer la sécurité et assurer la traçabilité des
opérations.

\subsubsection*{Question 11 (Moyen) : Le concept de « stratum » dans NTP}

\reponse
Le stratum représente le niveau hiérarchique d'un serveur NTP par
rapport à une référence de temps.
\begin{itemize}
    \item \textbf{Stratum 1 :} serveur directement synchronisé sur une
    source de temps de référence, comme une horloge GPS ou une horloge
    atomique.
    \item \textbf{Stratum 2 :} serveur synchronisé sur un ou plusieurs
    serveurs de stratum 1.
    \item \textbf{Stratum 3 :} serveur synchronisé sur un serveur de
    stratum 2.
\end{itemize}

\subsubsection*{Question 12 (Réflexion) : 3 exemples de services dépendant d'une synchronisation temporelle précise}

\reponse
\begin{enumerate}
    \item \textbf{SIEM et cybersécurité}

    Les outils SIEM collectent et analysent les logs provenant des
    serveurs, des pare-feu et des systèmes de détection d'intrusion. Une
    heure synchronisée permet de corréler les événements et de
    reconstituer la chronologie d'une cyberattaque.

    \item \textbf{Authentification et sécurité (Kerberos)}

    Le protocole Kerberos utilise des horodatages pour sécuriser
    l'authentification et limiter certaines attaques par rejeu. Une
    différence d'heure trop importante entre le client et le serveur
    peut provoquer des échecs d'authentification.

    \item \textbf{Transactions financières et applications métiers}

    Les systèmes financiers utilisent des horodatages pour assurer la
    traçabilité des opérations, faciliter les audits et reconstituer
    l'ordre des transactions. Une heure cohérente est également
    importante pour les journaux des applications métiers.
\end{enumerate}

\sources
\begin{itemize}
    \item IETF, RFC 2131 : Dynamic Host Configuration Protocol ;
    RFC 2132 : DHCP Options and BOOTP Vendor Extensions,
    \url{https://www.rfc-editor.org/rfc/rfc2131}.
    \item IETF, RFC 5905 : Network Time Protocol Version 4,
    \url{https://www.rfc-editor.org/rfc/rfc5905}.
\end{itemize}

% --------------------------------------------------
\subsection{Membre 3 : LDAP et RADIUS/TACACS+}
% --------------------------------------------------
\membre{\dotfill}{13 à 20}

\paragraph{A. LDAP (Lightweight Directory Access Protocol)}

\subsubsection*{Question 13 (Facile) : Qu'est-ce qu'un annuaire LDAP et à quoi sert-il ?}

\reponse
Un annuaire LDAP est une base de données spécialisée, organisée en
arbre, qui stocke des informations sur les ressources d'une
organisation : utilisateurs, groupes, ordinateurs, imprimantes,
services. LDAP est le protocole standard (TCP 389, ou 636 pour LDAPS)
qui permet de consulter et de modifier cet annuaire.

Il est optimisé pour la \textbf{lecture et la recherche rapides}, et non
pour les écritures fréquentes. Dans une organisation, il sert à :
\begin{itemize}
    \item \textbf{centraliser les identités} : un seul compte par
    utilisateur pour tous les services (messagerie, VPN, Wi-Fi, ENT) ;
    \item \textbf{authentifier} les utilisateurs (opération \textit{bind}) ;
    \item gérer les \textbf{groupes et les droits d'accès} ;
    \item fournir un \textbf{carnet d'adresses} commun (nom, e-mail,
    téléphone, service).
\end{itemize}

\exemple
À l'université, un étudiant utilise le même identifiant pour le Wi-Fi,
la plateforme e-learning et la messagerie : chaque service interroge le
même annuaire LDAP.

\sources
\begin{itemize}
    \item IETF, RFC 4511 : LDAP -- The Protocol,
    \url{https://www.rfc-editor.org/rfc/rfc4511}.
\end{itemize}

\subsubsection*{Question 14 (Moyen) : Structure hiérarchique d'un annuaire LDAP (DN, CN, OU, DC)}

\reponse
L'annuaire est organisé en arbre appelé \textbf{DIT} (Directory
Information Tree). Chaque nœud est une \textbf{entrée} composée
d'attributs.

\begin{longtable}{|L{2.2cm}|L{4.6cm}|L{7.8cm}|}
\hline
\textbf{Terme} & \textbf{Signification} & \textbf{Rôle} \\ \hline
\textbf{DC} & Domain Component & Composant du nom de domaine, placé à la
racine de l'arbre (ex. \texttt{dc=universite,dc=dz}). \\ \hline
\textbf{OU} & Organizational Unit & Conteneur qui regroupe des entrées
(département, service, type d'objet). \\ \hline
\textbf{CN} & Common Name & Nom d'un objet final : utilisateur, groupe,
machine. \\ \hline
\textbf{DN} & Distinguished Name & Identifiant unique et complet d'une
entrée : chemin depuis l'entrée jusqu'à la racine. \\ \hline
\end{longtable}

Le DN se lit de gauche à droite, de l'élément le plus précis vers la
racine. La partie la plus à gauche est le \textbf{RDN} (Relative
Distinguished Name).

\exemple
\begin{verbatim}
dc=universite,dc=dz
 |-- ou=Informatique
 |     |-- cn=Ahmed Benali
 |-- ou=Mathematiques
\end{verbatim}

\sources
\begin{itemize}
    \item IETF, RFC 4512 : LDAP Directory Information Models,
    \url{https://www.rfc-editor.org/rfc/rfc4512}.
\end{itemize}

\subsubsection*{Question 15 (Pratique) : DN complet de l'utilisateur « Ahmed Benali »}

\reponse
\begin{center}
\texttt{CN=Ahmed Benali,OU=Informatique,DC=universite,DC=dz}
\end{center}

\begin{itemize}
    \item \texttt{CN=Ahmed Benali} : l'utilisateur ;
    \item \texttt{OU=Informatique} : son département ;
    \item \texttt{DC=universite,DC=dz} : le domaine universite.dz.
\end{itemize}

\sources
\begin{itemize}
    \item IETF, RFC 4514 : LDAP String Representation of Distinguished
    Names, \url{https://www.rfc-editor.org/rfc/rfc4514}.
\end{itemize}

\subsubsection*{Question 16 (Avancé) : Comparaison LDAP et Active Directory}

\reponse
LDAP est un \textbf{protocole} d'accès à un annuaire. Active Directory
(AD) est un \textbf{service d'annuaire complet} de Microsoft, qui
utilise LDAP comme l'un de ses protocoles.

\textbf{Similitudes :} structure hiérarchique, même modèle de nommage
(DN, CN, OU, DC), recherche par requêtes LDAP, gestion des utilisateurs
et des groupes.

\begin{longtable}{|L{3.4cm}|L{5.6cm}|L{5.6cm}|}
\hline
\textbf{Critère} & \textbf{LDAP} & \textbf{Active Directory} \\ \hline
\endhead
Nature & Protocole ouvert (standard IETF) & Produit Microsoft
(Windows Server) \\ \hline
Implémentations & OpenLDAP, 389 Directory Server & AD DS \\ \hline
Authentification & Bind LDAP (simple ou SASL) & Kerberos (par défaut),
NTLM, LDAP \\ \hline
Services intégrés & Annuaire seulement & DNS, Kerberos, GPO, PKI,
réplication multi-maître \\ \hline
Gestion des postes & Non & Oui (stratégies de groupe, jonction au
domaine) \\ \hline
Plateforme & Multiplateforme & Surtout Windows \\ \hline
\end{longtable}

\exemple
Une entreprise sous Windows utilise AD pour gérer ses postes et ses
utilisateurs. Un serveur web Linux peut tout de même authentifier ces
utilisateurs en interrogeant AD via LDAP.

\sources
\begin{itemize}
    \item Microsoft Learn, Active Directory Domain Services Overview,
    \url{https://learn.microsoft.com/}.
\end{itemize}

\paragraph{B. RADIUS et TACACS+}

\subsubsection*{Question 17 (Facile) : Signification de AAA}

\reponse
\begin{itemize}
    \item \textbf{Authentication (Authentification) :} vérifier
    \textbf{qui} est l'utilisateur (identifiant et mot de passe,
    certificat, jeton).
    \item \textbf{Authorization (Autorisation) :} déterminer \textbf{ce
    qu'il a le droit de faire} (VLAN attribué, niveau de privilège,
    commandes permises).
    \item \textbf{Accounting (Traçabilité) :} enregistrer \textbf{ce
    qu'il a fait} (heure de connexion, durée, volume de données).
\end{itemize}

\exemple
Un employé se connecte au Wi-Fi : RADIUS vérifie son identité,
l'affecte au VLAN « Personnel », puis enregistre le début et la fin de
sa session.

\sources
\begin{itemize}
    \item IETF, RFC 2865 : RADIUS ; RFC 2866 : RADIUS Accounting,
    \url{https://www.rfc-editor.org/rfc/rfc2865}.
\end{itemize}

\subsubsection*{Question 18 (Moyen) : Flux d'authentification RADIUS}

\reponse
Trois acteurs interviennent : l'\textbf{utilisateur}, le \textbf{NAS}
(client RADIUS : switch, point d'accès, serveur VPN) et le
\textbf{serveur RADIUS}. Les échanges utilisent UDP 1812
(authentification) et 1813 (accounting).

\begin{enumerate}
    \item L'utilisateur fournit ses identifiants au NAS.
    \item Le NAS envoie un \textbf{Access-Request} au serveur RADIUS. Le
    mot de passe est chiffré avec un secret partagé entre le NAS et le
    serveur.
    \item Le serveur vérifie les identifiants (base locale, LDAP ou AD).
    \item Il répond par :
    \begin{itemize}
        \item \textbf{Access-Accept} : accès autorisé, avec des
        attributs (VLAN, ACL, durée de session) ;
        \item \textbf{Access-Reject} : accès refusé ;
        \item \textbf{Access-Challenge} : information supplémentaire
        demandée (code OTP, échange EAP).
    \end{itemize}
    \item Le NAS applique la décision.
    \item Le NAS envoie un \textbf{Accounting-Request (Start)}, puis
    \textbf{(Stop)} à la déconnexion.
\end{enumerate}

\sources
\begin{itemize}
    \item IETF, RFC 2865 : Remote Authentication Dial In User Service
    (RADIUS), \url{https://www.rfc-editor.org/rfc/rfc2865}.
\end{itemize}

\subsubsection*{Question 19 (Avancé) : RADIUS et TACACS+}

\reponse
\begin{longtable}{|L{3.4cm}|L{5.6cm}|L{5.6cm}|}
\hline
\textbf{Critère} & \textbf{RADIUS} & \textbf{TACACS+} \\ \hline
\endhead
Origine & Standard ouvert (IETF) & Développé par Cisco, publié en
RFC 8907 \\ \hline
Transport & UDP 1812/1813 & TCP 49 \\ \hline
Chiffrement & Seul le mot de passe est chiffré & Tout le corps du
paquet est chiffré \\ \hline
Séparation AAA & Authentification et autorisation combinées &
Authentification, autorisation et accounting séparés \\ \hline
Autorisation des commandes & Limitée & Commande par commande \\ \hline
Usage principal & Accès réseau des utilisateurs (Wi-Fi, 802.1X, VPN) &
Administration des équipements réseau \\ \hline
\end{longtable}

\textbf{Quand choisir l'un ou l'autre :}
\begin{itemize}
    \item \textbf{RADIUS} pour l'accès des utilisateurs au réseau : il
    est multi-constructeur et supporte 802.1X et EAP.
    \item \textbf{TACACS+} pour l'accès des administrateurs aux routeurs
    et switchs : chaque commande peut être autorisée et journalisée.
\end{itemize}
Les deux sont souvent utilisés ensemble dans une même entreprise.

\sources
\begin{itemize}
    \item IETF, RFC 8907 : The TACACS+ Protocol,
    \url{https://www.rfc-editor.org/rfc/rfc8907}.
\end{itemize}

\subsubsection*{Question 20 (Cas pratique) : Connexion SSH à un switch avec RADIUS}

\reponse
\begin{enumerate}
    \item L'administrateur lance \texttt{ssh admin@10.0.0.2}.
    \item Le switch établit la session SSH chiffrée et demande les
    identifiants.
    \item Le switch, configuré avec \texttt{aaa new-model} et un groupe
    RADIUS, envoie un \textbf{Access-Request} au serveur RADIUS.
    \item Le serveur vérifie l'utilisateur dans l'annuaire (AD/LDAP).
    \item Il renvoie un \textbf{Access-Accept} contenant le niveau de
    privilège (ex. \texttt{priv-lvl=15}) ou un \textbf{Access-Reject}.
    \item Le switch ouvre la session avec ce niveau de privilège.
    \item Un \textbf{Accounting-Start}, puis un \textbf{Accounting-Stop},
    sont envoyés pour tracer la session.
    \item Si le serveur RADIUS est injoignable, le switch peut utiliser
    un compte local de secours (\texttt{local}).
\end{enumerate}

\exemple
\begin{verbatim}
aaa new-model
radius server RAD1
 address ipv4 10.0.0.10 auth-port 1812 acct-port 1813
 key S3cret
aaa authentication login default group radius local
aaa authorization exec default group radius local
line vty 0 4
 transport input ssh
\end{verbatim}

\sources
\begin{itemize}
    \item Cisco, Configuring RADIUS, documentation IOS,
    \url{https://www.cisco.com/}.
\end{itemize}

% --------------------------------------------------
\subsection{Membre 4 : Kerberos et SNMP}
% --------------------------------------------------
\membre{\dotfill}{21 à 27}

\paragraph{C. Kerberos}

\subsubsection*{Question 21 (Moyen) : Le principe du « ticket » dans Kerberos}

\reponse
Un \textbf{ticket} est un justificatif chiffré, délivré par un tiers de
confiance (le KDC), qui prouve l'identité d'un utilisateur auprès d'un
service. Il contient l'identité de l'utilisateur, une clé de session et
une durée de validité (souvent 10 heures).

Son rôle : l'utilisateur s'authentifie \textbf{une seule fois} avec son
mot de passe, puis présente des tickets aux services. Le mot de passe
ne circule jamais sur le réseau. C'est le principe du \textbf{Single
Sign-On (SSO)}.

\exemple
Le matin, un employé ouvre sa session Windows. Il accède ensuite aux
partages de fichiers, à l'intranet et à l'imprimante sans retaper son
mot de passe : ses tickets Kerberos sont présentés automatiquement.

\sources
\begin{itemize}
    \item IETF, RFC 4120 : The Kerberos Network Authentication Service
    (V5), \url{https://www.rfc-editor.org/rfc/rfc4120}.
\end{itemize}

\subsubsection*{Question 22 (Avancé) : KDC, TGT et TGS}

\reponse
\begin{itemize}
    \item \textbf{KDC (Key Distribution Center) :} serveur central de
    confiance, qui connaît les clés secrètes de tous les utilisateurs et
    services. Il comprend l'\textbf{AS} (Authentication Service) et le
    \textbf{TGS}. Dans Windows, c'est le contrôleur de domaine.
    \item \textbf{TGT (Ticket Granting Ticket) :} ticket délivré par
    l'AS après l'authentification initiale. Il sert à demander d'autres
    tickets sans redonner son mot de passe.
    \item \textbf{TGS (Ticket Granting Service) :} service du KDC qui
    délivre les \textbf{tickets de service} en échange du TGT.
\end{itemize}

\textbf{Interaction (port 88) :}
\begin{enumerate}
    \item \textbf{AS-REQ / AS-REP :} le client s'authentifie auprès de
    l'AS et reçoit un TGT et une clé de session.
    \item \textbf{TGS-REQ / TGS-REP :} le client présente le TGT au TGS
    et reçoit un ticket pour le service voulu.
    \item \textbf{AP-REQ / AP-REP :} le client présente ce ticket au
    serveur d'application, qui le déchiffre avec sa propre clé et
    accorde l'accès.
\end{enumerate}

\sources
\begin{itemize}
    \item IETF, RFC 4120, \url{https://www.rfc-editor.org/rfc/rfc4120}.
\end{itemize}

\subsubsection*{Question 23 (Réflexion) : Pourquoi Kerberos est-il plus sécurisé qu'un mot de passe simple ?}

\reponse
\begin{itemize}
    \item \textbf{Le mot de passe ne circule jamais} sur le réseau, ni
    en clair ni chiffré. Il sert seulement à dériver une clé locale.
    \item \textbf{Authentification mutuelle :} le client peut vérifier
    l'identité du serveur, ce qui limite les faux serveurs.
    \item \textbf{Tickets limités dans le temps :} un ticket volé expire
    rapidement.
    \item \textbf{Protection contre le rejeu :} les authentificateurs
    contiennent un horodatage, d'où la tolérance d'horloge de
    5~minutes par défaut (et l'importance de NTP).
    \item \textbf{SSO :} moins de saisies de mot de passe, donc moins
    d'exposition.
\end{itemize}

\textbf{Limites :} le KDC est un point critique, et des attaques
existent (Pass-the-Ticket, Golden Ticket, Kerberoasting).

\sources
\begin{itemize}
    \item IETF, RFC 4120, \url{https://www.rfc-editor.org/rfc/rfc4120}.
\end{itemize}

\paragraph{A. SNMP (Simple Network Management Protocol)}

\subsubsection*{Question 24 (Facile) : Les trois composants d'une architecture SNMP}

\reponse
\begin{enumerate}
    \item \textbf{Le manager (NMS, Network Management Station) :} poste
    de supervision qui interroge les équipements et reçoit les alertes
    (Zabbix, LibreNMS, PRTG).
    \item \textbf{L'agent SNMP :} logiciel présent sur chaque équipement
    supervisé (routeur, switch, serveur), qui répond aux requêtes et
    envoie des alertes.
    \item \textbf{La MIB :} base d'informations qui décrit les objets
    consultables sur l'agent.
\end{enumerate}
Ports utilisés : UDP 161 (requêtes vers l'agent), UDP 162 (traps vers
le manager).

\sources
\begin{itemize}
    \item IETF, RFC 3411 : An Architecture for Describing SNMP
    Management Frameworks, \url{https://www.rfc-editor.org/rfc/rfc3411}.
\end{itemize}

\subsubsection*{Question 25 (Moyen) : SNMPv1, SNMPv2c et SNMPv3}

\reponse
\begin{longtable}{|L{2.4cm}|L{6cm}|L{6.2cm}|}
\hline
\textbf{Version} & \textbf{Caractéristiques} & \textbf{Sécurité} \\ \hline
\endhead
SNMPv1 & Première version. Opérations GET, GETNEXT, SET, TRAP. Compteurs
32 bits. & Chaîne de communauté envoyée \textbf{en clair}. \\ \hline
SNMPv2c & Ajoute \textbf{GETBULK} (lecture en masse), \textbf{INFORM}
(trap acquitté), compteurs 64 bits, meilleurs messages d'erreur. &
Toujours la communauté en clair. \\ \hline
SNMPv3 & Même fonctionnalités que v2c, avec un vrai modèle de sécurité
(USM et VACM). & \textbf{Authentification} (SHA), \textbf{chiffrement}
(AES), contrôle d'accès par vue. \\ \hline
\end{longtable}

\textbf{Pourquoi SNMPv3 est recommandé :} avec v1 et v2c, un attaquant
qui capture le trafic récupère la communauté et peut lire, voire
modifier (SET), la configuration des équipements. SNMPv3 garantit
l'authentification, l'intégrité et la confidentialité, avec trois
niveaux : \textit{noAuthNoPriv}, \textit{authNoPriv} et
\textit{authPriv} (recommandé).

\sources
\begin{itemize}
    \item IETF, RFC 3414 : User-based Security Model (USM) for SNMPv3,
    \url{https://www.rfc-editor.org/rfc/rfc3414}.
\end{itemize}

\subsubsection*{Question 26 (Pratique) : Qu'est-ce qu'une MIB ?}

\reponse
La \textbf{MIB (Management Information Base)} est une base de données
hiérarchique qui décrit tous les objets qu'un agent SNMP peut exposer.
Chaque objet est identifié par un \textbf{OID} (Object Identifier),
une suite de nombres qui représente son chemin dans l'arbre (ex.
\texttt{1.3.6.1.2.1.1.1}). Il existe des MIB standards (MIB-II) et des
MIB propriétaires des constructeurs.

\exemple
\begin{longtable}{|L{4cm}|L{4.6cm}|L{6cm}|}
\hline
\textbf{Objet} & \textbf{OID} & \textbf{Information} \\ \hline
sysDescr & 1.3.6.1.2.1.1.1 & Description de l'équipement (modèle,
version) \\ \hline
sysUpTime & 1.3.6.1.2.1.1.3 & Temps écoulé depuis le démarrage \\ \hline
sysName & 1.3.6.1.2.1.1.5 & Nom de l'équipement \\ \hline
ifOperStatus & 1.3.6.1.2.1.2.2.1.8 & État d'une interface (up/down) \\ \hline
ifInOctets & 1.3.6.1.2.1.2.2.1.10 & Octets reçus sur une interface \\ \hline
\end{longtable}
On y trouve aussi l'utilisation CPU et mémoire, la température, l'état
des ventilateurs et les tables de routage.

\sources
\begin{itemize}
    \item IETF, RFC 1213 : MIB-II,
    \url{https://www.rfc-editor.org/rfc/rfc1213}.
\end{itemize}

\subsubsection*{Question 27 (Avancé) : Opérations GET, SET et TRAP}

\reponse
\begin{longtable}{|L{1.8cm}|L{3.2cm}|L{4.8cm}|L{4.8cm}|}
\hline
\textbf{Opération} & \textbf{Sens} & \textbf{Rôle} & \textbf{Contexte} \\ \hline
\endhead
GET & Manager \fl{} agent & Lire la valeur d'un objet & Supervision
périodique (polling) : débit, CPU, état des interfaces \\ \hline
SET & Manager \fl{} agent & Modifier la valeur d'un objet & Action à
distance : désactiver un port, changer un paramètre \\ \hline
TRAP & Agent \fl{} manager & Notification spontanée, sans accusé de
réception & Alerte immédiate : lien coupé, redémarrage, seuil dépassé
\\ \hline
\end{longtable}

GET fonctionne en mode « pull » (le manager demande), TRAP en mode
« push » (l'agent prévient). Les deux sont complémentaires. En SNMPv2c
et v3, l'\textbf{INFORM} est un trap avec accusé de réception.

\sources
\begin{itemize}
    \item IETF, RFC 3416 : Protocol Operations for SNMP,
    \url{https://www.rfc-editor.org/rfc/rfc3416}.
\end{itemize}

% --------------------------------------------------
\subsection{Membre 5 : Syslog et NetFlow}
% --------------------------------------------------
\membre{\dotfill}{28 à 33}

\paragraph{B. Syslog}

\subsubsection*{Question 28 (Facile) : Objectif d'un serveur Syslog centralisé}

\reponse
Un serveur Syslog centralisé \textbf{collecte en un seul endroit les
journaux (logs)} de tous les équipements : routeurs, switchs, pare-feu,
serveurs. Les messages sont envoyés en UDP 514 (ou TCP, et TLS sur le
port 6514 pour la sécurité).

Il permet de conserver les logs même si un équipement redémarre ou est
compromis, de les consulter depuis une seule interface et de les
transmettre à un SIEM.

\sources
\begin{itemize}
    \item IETF, RFC 5424 : The Syslog Protocol,
    \url{https://www.rfc-editor.org/rfc/rfc5424}.
\end{itemize}

\subsubsection*{Question 29 (Moyen) : Les niveaux de sévérité Syslog (0 à 7)}

\reponse
Plus le chiffre est bas, plus le message est grave.

\begin{longtable}{|L{0.9cm}|L{3cm}|L{10.6cm}|}
\hline
\textbf{N°} & \textbf{Niveau} & \textbf{Exemple} \\ \hline
\endhead
0 & Emergency & Le système est inutilisable (panne totale, kernel
panic). \\ \hline
1 & Alert & Action immédiate requise (base de données corrompue,
surchauffe). \\ \hline
2 & Critical & Condition critique (défaillance d'un disque ou d'une
alimentation). \\ \hline
3 & Error & Erreur (interface en échec, échec d'un service). \\ \hline
4 & Warning & Avertissement (CPU à 85~\%, espace disque faible). \\ \hline
5 & Notice & Événement normal mais important (interface qui passe up
ou down, changement de configuration). \\ \hline
6 & Informational & Information (utilisateur connecté, ACL appliquée).
\\ \hline
7 & Debug & Messages de débogage (détails d'un protocole). \\ \hline
\end{longtable}

\sources
\begin{itemize}
    \item IETF, RFC 5424, \url{https://www.rfc-editor.org/rfc/rfc5424}.
\end{itemize}

\subsubsection*{Question 30 (Pratique) : Pourquoi centraliser les logs ?}

\reponse
\textbf{En sécurité :}
\begin{itemize}
    \item un attaquant qui compromet un équipement ne peut pas effacer
    les logs déjà envoyés au serveur central ;
    \item la \textbf{corrélation} d'événements entre plusieurs
    équipements permet de détecter une attaque (force brute, scan,
    mouvement latéral) ;
    \item conservation pour l'\textbf{audit}, l'investigation
    \textit{forensic} et la conformité réglementaire.
\end{itemize}

\textbf{En dépannage :}
\begin{itemize}
    \item vue unique au lieu de se connecter à chaque équipement ;
    \item reconstitution de la chronologie d'un incident (à condition
    que les horloges soient synchronisées avec NTP) ;
    \item recherche, filtrage et alertes automatiques.
\end{itemize}

\exemple
Dix tentatives SSH échouées sur plusieurs switchs en une minute, depuis
la même IP : le SIEM qui reçoit les logs centralisés détecte une
attaque par force brute et déclenche une alerte.

\sources
\begin{itemize}
    \item NIST, SP 800-92 : Guide to Computer Security Log Management,
    \url{https://csrc.nist.gov/}.
\end{itemize}

\paragraph{C. NetFlow}

\textit{Remarque : vérifier dans le sujet original le libellé
exact de la question 31, dont le texte est incomplet dans
la version fournie.}

\subsubsection*{Question 31 : NetFlow et informations collectées sur le trafic réseau}

\reponse
NetFlow est une technologie (créée par Cisco, standardisée ensuite sous
le nom \textbf{IPFIX}) qui permet aux routeurs et switchs d'exporter
des \textbf{statistiques sur les flux} qui les traversent vers un
\textbf{collecteur}. Un flux est un ensemble de paquets qui partagent
les mêmes caractéristiques.

Il ne capture pas le contenu des paquets, mais les
\textbf{métadonnées} : adresses IP et ports source et destination,
protocole, nombre de paquets et d'octets, heures de début et de fin,
interfaces d'entrée et de sortie, flags TCP, ToS/DSCP, numéros d'AS.
Il répond aux questions « qui parle à qui, avec quel protocole, et
combien ? ».

\sources
\begin{itemize}
    \item IETF, RFC 3954 : Cisco Systems NetFlow Services Export
    Version 9 ; RFC 7011 : IPFIX,
    \url{https://www.rfc-editor.org/rfc/rfc7011}.
\end{itemize}

\subsubsection*{Question 32 (Avancé) : Les 7 champs clés d'un flow NetFlow}

\reponse
Un flux est identifié par 7 champs. Deux paquets appartiennent au même
flux si ces 7 valeurs sont identiques :
\begin{enumerate}
    \item Adresse IP source ;
    \item Adresse IP destination ;
    \item Port source ;
    \item Port destination ;
    \item Protocole de couche 3 (TCP, UDP, ICMP) ;
    \item Octet ToS (Type of Service / DSCP) ;
    \item Interface d'entrée (ifIndex).
\end{enumerate}

\sources
\begin{itemize}
    \item Cisco, Introduction to Cisco IOS NetFlow,
    \url{https://www.cisco.com/}.
\end{itemize}

\subsubsection*{Question 33 (Cas pratique) : Utilisation de NetFlow}

\reponse
\textbf{Détecter une attaque DDoS :}
\begin{itemize}
    \item hausse brutale du nombre de flux ou de paquets vers une même
    IP de destination ;
    \item très grand nombre d'IP sources différentes vers une seule
    cible ;
    \item flux de petits paquets (SYN flood) ou trafic UDP massif vers
    des ports d'amplification (53, 123, 1900) ;
    \item alerte sur seuil, puis mitigation (RTBH, filtrage).
\end{itemize}

\textbf{Identifier les applications qui consomment le plus de bande
passante :}
\begin{itemize}
    \item classement « Top N » par octets, selon les ports et protocoles
    (443, 1935, etc.) ou par IP/utilisateur ;
    \item suivi historique pour la planification de capacité et le
    réglage de la QoS.
\end{itemize}

\textbf{Analyser les patterns de trafic anormaux :}
\begin{itemize}
    \item établir une \textbf{référence (baseline)} du trafic normal,
    puis détecter les écarts ;
    \item scan de ports (une IP vers beaucoup de ports), communications
    vers des pays ou IP suspectes, exfiltration (gros volumes sortants la
    nuit), trafic sur des ports inhabituels.
\end{itemize}

\exemple
Outils : ntopng, Elastiflow, SolarWinds NTA, ou un SIEM qui reçoit les
flux IPFIX.

\sources
\begin{itemize}
    \item IETF, RFC 7011, \url{https://www.rfc-editor.org/rfc/rfc7011}.
\end{itemize}

% --------------------------------------------------
\subsection{Membre 6 : HTTP/HTTPS et serveurs Web}
% --------------------------------------------------
\membre{\dotfill}{42 à 45}

\paragraph{A. HTTP/HTTPS et serveurs Web}

\subsubsection*{Question 42 (Facile) : Différence entre HTTP et HTTPS}

\reponse
\begin{itemize}
    \item \textbf{HTTP} (port 80) transmet les données \textbf{en clair}
    : un attaquant sur le chemin peut lire et modifier les échanges.
    \item \textbf{HTTPS} (port 443) est HTTP encapsulé dans
    \textbf{TLS}. Il apporte la \textbf{confidentialité} (chiffrement),
    l'\textbf{intégrité} (données non modifiées) et
    l'\textbf{authentification} du serveur (certificat).
\end{itemize}

\textbf{Pourquoi HTTPS est devenu la norme :}
\begin{itemize}
    \item protection des mots de passe, paiements et données
    personnelles ;
    \item les navigateurs affichent « Non sécurisé » pour les sites HTTP
    et les moteurs de recherche favorisent HTTPS ;
    \item certificats gratuits et automatisés (Let's Encrypt) ;
    \item HTTP/2 et HTTP/3 ne sont utilisés par les navigateurs qu'avec
    TLS ;
    \item exigences réglementaires (protection des données, PCI DSS).
\end{itemize}

\sources
\begin{itemize}
    \item IETF, RFC 9110 : HTTP Semantics,
    \url{https://www.rfc-editor.org/rfc/rfc9110}.
\end{itemize}

\subsubsection*{Question 43 (Moyen) : Les méthodes HTTP principales}

\reponse
\begin{longtable}{|L{1.8cm}|L{6cm}|L{2.4cm}|L{4.4cm}|}
\hline
\textbf{Méthode} & \textbf{Rôle} & \textbf{Idempotente} &
\textbf{Exemple} \\ \hline
\endhead
GET & Lire une ressource, sans la modifier & Oui &
\texttt{GET /etudiants/12} \\ \hline
POST & Créer une ressource ou envoyer des données à traiter & Non &
\texttt{POST /etudiants} \\ \hline
PUT & Remplacer entièrement une ressource & Oui &
\texttt{PUT /etudiants/12} \\ \hline
PATCH & Modifier partiellement une ressource & Pas forcément &
\texttt{PATCH /etudiants/12} (e-mail seul) \\ \hline
DELETE & Supprimer une ressource & Oui & \texttt{DELETE /etudiants/12}
\\ \hline
\end{longtable}
Une méthode \textbf{idempotente} donne le même résultat si on la répète
plusieurs fois.

\sources
\begin{itemize}
    \item IETF, RFC 9110 ; RFC 5789 : PATCH Method for HTTP,
    \url{https://www.rfc-editor.org/rfc/rfc5789}.
\end{itemize}

\subsubsection*{Question 44 (Avancé) : HTTP/2 et HTTP/3}

\reponse
HTTP/1.1 envoie une requête à la fois par connexion TCP, avec des
en-têtes texte répétés. Il souffre du \textbf{blocage en tête de file}
(\textit{head-of-line blocking}), et les navigateurs ouvrent plusieurs
connexions pour compenser.

\begin{itemize}
    \item \textbf{Multiplexage (HTTP/2) :} plusieurs requêtes et
    réponses circulent en parallèle sur \textbf{une seule connexion},
    sous forme de flux (\textit{streams}) découpés en trames binaires.
    \item \textbf{Server Push (HTTP/2) :} le serveur envoie des
    ressources avant qu'elles soient demandées (CSS, JS). En pratique,
    il a été abandonné par les principaux navigateurs, au profit des
    \textit{preload hints} (réponse 103 Early Hints).
    \item \textbf{Compression des en-têtes :} \textbf{HPACK} en HTTP/2,
    \textbf{QPACK} en HTTP/3. Les en-têtes répétés ne sont envoyés
    qu'une fois.
    \item \textbf{QUIC (HTTP/3) :} HTTP/3 remplace TCP par QUIC, sur
    UDP. QUIC intègre TLS 1.3, établit la connexion plus vite (1-RTT,
    voire 0-RTT), supprime le blocage en tête de file au niveau
    transport (la perte d'un paquet ne bloque qu'un seul flux) et permet
    la \textbf{migration de connexion} (passage du Wi-Fi à la 4G sans
    coupure).
\end{itemize}

\sources
\begin{itemize}
    \item IETF, RFC 9113 : HTTP/2 ; RFC 9114 : HTTP/3 ; RFC 9000 :
    QUIC, \url{https://www.rfc-editor.org/rfc/rfc9114}.
\end{itemize}

\subsubsection*{Question 45 (Pratique) : Les codes de statut HTTP}

\reponse
\begin{longtable}{|L{4.4cm}|L{10.2cm}|}
\hline
\textbf{Code} & \textbf{Contexte} \\ \hline
\endhead
200 OK & Requête réussie, la ressource est renvoyée. \\ \hline
301 Moved Permanently & La ressource a changé définitivement d'adresse
(ex. redirection de HTTP vers HTTPS). \\ \hline
304 Not Modified & La ressource n'a pas changé depuis la version en
cache du client : rien n'est renvoyé. \\ \hline
401 Unauthorized & Authentification requise ou identifiants invalides.
\\ \hline
403 Forbidden & Le client est identifié, mais n'a pas le droit
d'accéder à la ressource. \\ \hline
404 Not Found & La ressource n'existe pas. \\ \hline
500 Internal Server Error & Erreur dans l'application côté serveur (bug,
exception). \\ \hline
502 Bad Gateway & Un reverse proxy ou une passerelle a reçu une réponse
invalide du serveur en amont (serveur d'application arrêté). \\ \hline
503 Service Unavailable & Serveur surchargé ou en maintenance,
indisponibilité temporaire. \\ \hline
\end{longtable}
Familles : 2xx succès, 3xx redirection, 4xx erreur du client, 5xx
erreur du serveur.

\sources
\begin{itemize}
    \item IETF, RFC 9110, section 15,
    \url{https://www.rfc-editor.org/rfc/rfc9110}.
\end{itemize}

% --------------------------------------------------
\subsection{Membre 7 : Proxy, reverse proxy et CDN}
% --------------------------------------------------
\membre{\dotfill}{46 à 51}

\paragraph{B. Proxy et Reverse Proxy}

\subsubsection*{Question 46 (Moyen) : Forward proxy et reverse proxy}

\reponse
\begin{itemize}
    \item \textbf{Forward proxy :} placé \textbf{devant les clients}. Il
    reçoit leurs requêtes et les transmet à Internet en leur nom. Il
    représente les clients.

    \textit{Exemples :} filtrage web en entreprise (Squid), contrôle
    d'accès, journalisation, cache, contournement de restrictions.

    \item \textbf{Reverse proxy :} placé \textbf{devant les serveurs}.
    Il reçoit les requêtes d'Internet et les transmet aux serveurs
    internes. Il représente les serveurs.

    \textit{Exemples :} Nginx, HAProxy ou Traefik devant des serveurs
    web : terminaison TLS, répartition de charge, cache, WAF.
\end{itemize}

\sources
\begin{itemize}
    \item Nginx, documentation Reverse Proxy,
    \url{https://docs.nginx.com/}.
\end{itemize}

\subsubsection*{Question 47 (Pratique) : Avantages d'un reverse proxy}

\reponse
\begin{enumerate}
    \item \textbf{Répartition de charge} entre plusieurs serveurs.
    \item \textbf{Terminaison TLS} centralisée : un seul endroit pour
    gérer les certificats, et moins de charge sur les serveurs.
    \item \textbf{Cache} des contenus fréquents.
    \item \textbf{Compression} (gzip, Brotli) des réponses.
    \item \textbf{Masquage} de l'architecture interne : les clients ne
    voient jamais l'adresse réelle des serveurs.
    \item \textbf{Sécurité} : WAF, limitation de débit, filtrage.
    \item \textbf{Haute disponibilité} : health checks et bascule.
    \item \textbf{Routage} selon l'URL ou le nom de domaine vers
    différentes applications.
\end{enumerate}

\sources
\begin{itemize}
    \item HAProxy Technologies, documentation,
    \url{https://www.haproxy.com/documentation/}.
\end{itemize}

\subsubsection*{Question 48 (Avancé) : Sécurité, performances et disponibilité}

\reponse
\textbf{Sécurité :}
\begin{itemize}
    \item il cache les IP et la technologie des serveurs (suppression de
    l'en-tête \texttt{Server}) ;
    \item WAF (ModSecurity) contre les injections SQL et le XSS ;
    \item limitation de débit (\textit{rate limiting}) contre la force
    brute et les petites attaques DDoS ;
    \item TLS imposé partout et redirection HTTP \fl{} HTTPS ;
    \item filtrage par IP, en-têtes ou pays.
\end{itemize}

\textbf{Performances :}
\begin{itemize}
    \item cache des réponses (images, pages), qui n'atteignent plus les
    serveurs ;
    \item compression gzip ou Brotli ;
    \item réutilisation des connexions (\textit{keep-alive}) vers les
    serveurs ;
    \item HTTP/2 et HTTP/3 côté clients, même si les serveurs ne les
    supportent pas.
\end{itemize}

\textbf{Disponibilité :}
\begin{itemize}
    \item load balancing entre plusieurs serveurs ;
    \item health checks : un serveur en panne est retiré automatiquement ;
    \item failover vers un serveur de secours ;
    \item maintenance sans coupure (retrait progressif d'un serveur).
\end{itemize}

\sources
\begin{itemize}
    \item OWASP, ModSecurity Core Rule Set, \url{https://owasp.org/}.
\end{itemize}

\paragraph{C. CDN (Content Delivery Network)}

\subsubsection*{Question 49 (Facile) : Qu'est-ce qu'un CDN ?}

\reponse
Un CDN est un \textbf{réseau de serveurs de cache} (\textit{edge
servers}) répartis dans le monde. Il stocke une copie des contenus
d'un site au plus près des utilisateurs.

\textbf{Principe :}
\begin{enumerate}
    \item l'utilisateur demande une ressource ;
    \item le DNS ou l'Anycast le dirige vers le serveur edge le plus
    proche ;
    \item si la ressource est en cache (\textit{cache hit}), elle est
    servie directement ;
    \item sinon (\textit{cache miss}), l'edge la récupère sur le serveur
    d'origine, la sert et la garde en cache.
\end{enumerate}
Résultat : latence réduite, moins de charge sur l'origine, absorption
des pics de trafic et d'une partie des attaques DDoS.

\exemple
Cloudflare, Akamai, Amazon CloudFront.

\sources
\begin{itemize}
    \item Cloudflare Learning Center, What is a CDN?,
    \url{https://www.cloudflare.com/learning/}.
\end{itemize}

\subsubsection*{Question 50 (Moyen) : Contenus qui bénéficient le plus d'un CDN}

\reponse
\begin{itemize}
    \item \textbf{Contenu statique} (images, CSS, JavaScript, polices) :
    identique pour tous les utilisateurs et rarement modifié, donc
    facile à mettre en cache longtemps.
    \item \textbf{Vidéo et streaming} (HLS, DASH) : gros volumes, très
    sensibles à la latence et au débit.
    \item \textbf{Téléchargements de fichiers} (mises à jour logicielles,
    jeux, ISO) : volume énorme et demandes simultanées.
    \item \textbf{Pages web peu changeantes} (blogs, documentation).
\end{itemize}
Les contenus \textbf{dynamiques et personnalisés} (panier, compte
utilisateur) en profitent moins pour le cache, mais gagnent quand même
en accélération réseau et en protection DDoS.

\sources
\begin{itemize}
    \item Akamai, What Is a CDN?, \url{https://www.akamai.com/}.
\end{itemize}

\subsubsection*{Question 51 (Avancé) : Stratégies de cache dans un CDN}

\reponse
\begin{itemize}
    \item \textbf{En-têtes Cache-Control :} le serveur d'origine indique
    comment cacher. \texttt{max-age} (durée côté navigateur),
    \texttt{s-maxage} (durée côté CDN), \texttt{public} ou
    \texttt{private}, \texttt{no-store} (ne jamais cacher),
    \texttt{no-cache} (revalider à chaque fois).
    \item \textbf{TTL (Time To Live) :} durée pendant laquelle une copie
    reste valide. Long pour le statique (un an avec des noms de fichiers
    versionnés), court ou nul pour le dynamique. À expiration, l'edge
    revalide avec \texttt{ETag} ou \texttt{Last-Modified} (réponse 304
    si rien n'a changé).
    \item \textbf{Edge caching et Origin shield :} l'edge cache est au
    plus près des utilisateurs. L'\textbf{origin shield} est une couche
    de cache intermédiaire entre les edges et l'origine : tous les cache
    miss passent par lui, ce qui protège l'origine d'un trop grand
    nombre de requêtes.
    \item \textbf{Invalidation de cache :} supprimer une copie avant la
    fin de son TTL, par \textbf{purge} (par URL, par tag ou complète) ou
    par \textbf{versionnement} des fichiers (\texttt{style.v2.css}), qui
    est la méthode la plus fiable.
\end{itemize}

\exemple
\texttt{Cache-Control: public, max-age=31536000, immutable} pour
\texttt{app.3f2a1.js}, et \texttt{Cache-Control: no-store} pour la page
de profil d'un utilisateur.

\sources
\begin{itemize}
    \item IETF, RFC 9111 : HTTP Caching,
    \url{https://www.rfc-editor.org/rfc/rfc9111}.
\end{itemize}

% --------------------------------------------------
\subsection{Membre 8 : Load balancing et DNS Web}
% --------------------------------------------------
\membre{\dotfill}{52 à 57}

\paragraph{D. Load Balancing Web}

\subsubsection*{Question 52 (Moyen) : Algorithmes de répartition de charge}

\reponse
\begin{longtable}{|L{3cm}|L{5.8cm}|L{5.8cm}|}
\hline
\textbf{Algorithme} & \textbf{Principe} & \textbf{Scénario adapté} \\ \hline
\endhead
Round Robin & Les requêtes sont distribuées à tour de rôle. & Serveurs
identiques, requêtes de durée similaire. \\ \hline
Least Connections & La requête va au serveur qui a le moins de
connexions actives. & Requêtes de durée variable (sessions longues,
WebSocket, base de données). \\ \hline
IP Hash & L'IP du client détermine toujours le même serveur. & Besoin
de persistance de session sans cookie. \\ \hline
Weighted Round Robin & Tour de rôle pondéré selon la capacité de chaque
serveur. & Serveurs de puissances différentes (un serveur deux fois
plus puissant reçoit deux fois plus de requêtes). \\ \hline
\end{longtable}

\sources
\begin{itemize}
    \item Nginx, HTTP Load Balancing, \url{https://docs.nginx.com/}.
\end{itemize}

\subsubsection*{Question 53 (Avancé) : Load balancing de couche 4 et de couche 7}

\reponse
\begin{itemize}
    \item \textbf{Couche 4 (Transport) :} décide uniquement à partir des
    IP, des ports et du protocole (TCP/UDP). Il ne lit pas le contenu.
    Il est très rapide et peu gourmand, et fonctionne pour tout
    protocole.

    \textit{Décisions possibles :} tout le trafic TCP 443 vers le pool
    A, le trafic UDP 53 vers le pool DNS, répartition par IP source.

    \item \textbf{Couche 7 (Application) :} lit le contenu HTTP (URL,
    en-têtes, cookies, nom d'hôte). Il est plus intelligent mais plus
    coûteux, et doit déchiffrer TLS.

    \textit{Décisions possibles :} \texttt{/api/*} vers les serveurs
    d'API, \texttt{/images/*} vers les serveurs statiques,
    \texttt{shop.exemple.dz} vers l'application boutique, utilisateurs
    mobiles (User-Agent) vers une version dédiée, persistance par
    cookie, tests A/B.
\end{itemize}

\exemple
Couche 4 : AWS Network Load Balancer, LVS. Couche 7 : HAProxy en mode
HTTP, Nginx, AWS Application Load Balancer.

\sources
\begin{itemize}
    \item HAProxy, Layer 4 vs Layer 7 Load Balancing,
    \url{https://www.haproxy.com/}.
\end{itemize}

\subsubsection*{Question 54 (Pratique) : Architecture web complète}

\reponse
\begin{center}
\small
Utilisateurs \fl{} DNS (GeoDNS) \fl{} \textbf{CDN} \fl{} \textbf{WAF +
Reverse proxy / LB couche 7} (en paire HA)\\
\fl{} \textbf{Pool API} / \textbf{Pool statique} / \textbf{Serveurs de
streaming}\\
\fl{} \textbf{Cache Redis} \fl{} \textbf{Base de données} (primaire +
réplique)
\end{center}

\textbf{Justification de chaque choix :}
\begin{itemize}
    \item \textbf{CDN :} sert les pages statiques et les segments vidéo
    au plus près des utilisateurs, absorbe les pics et les DDoS.
    \item \textbf{Reverse proxy / load balancer couche 7} (HAProxy ou
    Nginx, en paire active/passive avec VRRP) : terminaison TLS, WAF,
    routage par URL (\texttt{/api} \fl{} pool API, \texttt{/video} \fl{}
    streaming), health checks.
    \item \textbf{Serveurs d'application de l'API REST :} plusieurs
    instances sans état (stateless), ajoutées selon la charge, avec
    l'algorithme least connections.
    \item \textbf{Streaming vidéo :} serveurs dédiés (HLS/DASH) ou
    stockage objet derrière le CDN, car le trafic est très volumineux.
    \item \textbf{Cache (Redis ou Memcached) :} stocke les résultats
    fréquents de l'API et les sessions, ce qui soulage la base.
    \item \textbf{Base de données :} un serveur primaire pour les
    écritures et des répliques en lecture, avec sauvegardes, dans un
    réseau privé non exposé.
\end{itemize}

\sources
\begin{itemize}
    \item AWS, Well-Architected Framework,
    \url{https://aws.amazon.com/architecture/well-architected/}.
\end{itemize}

\paragraph{E. DNS pour les Services Web}

\subsubsection*{Question 55 (Moyen) : DNS Round Robin et ses limites}

\reponse
Le \textbf{DNS Round Robin} consiste à associer \textbf{plusieurs
enregistrements A} au même nom. Le serveur DNS renvoie les adresses
dans un ordre qui change à chaque réponse, ce qui répartit les clients
entre les serveurs.

\textbf{Limites :}
\begin{itemize}
    \item \textbf{pas de vérification de santé} : une IP en panne
    continue d'être distribuée ;
    \item \textbf{cache DNS} des résolveurs et des clients : la
    répartition est inégale et les changements sont lents ;
    \item pas de prise en compte de la charge ni de la capacité des
    serveurs ;
    \item pas de persistance de session ;
    \item le client garde l'IP reçue, même si elle devient lente.
\end{itemize}

\sources
\begin{itemize}
    \item IETF, RFC 1794 : DNS Support for Load Balancing,
    \url{https://www.rfc-editor.org/rfc/rfc1794}.
\end{itemize}

\subsubsection*{Question 56 (Avancé) : GeoDNS, DNS Failover, DNS Load Balancing et Anycast}

\reponse
\begin{itemize}
    \item \textbf{GeoDNS :} le DNS répond avec une IP différente selon la
    localisation du client. Il dirige l'utilisateur vers le datacenter
    le plus proche (latence, conformité légale).
    \item \textbf{DNS Failover :} des sondes surveillent les serveurs.
    En cas de panne, le DNS retire l'IP défaillante et renvoie celle du
    site de secours. Un TTL court (30 à 60 s) accélère la bascule.
    \item \textbf{DNS Load Balancing :} évolution du Round Robin avec
    health checks, pondération et prise en compte de la latence
    (Route 53, Cloudflare Load Balancing).
    \item \textbf{Anycast :} la même IP est annoncée en BGP depuis
    plusieurs sites. Le routage Internet envoie chaque client vers le
    site le plus proche. Il est utilisé par les serveurs DNS racine et
    les CDN, et résiste bien aux DDoS.
\end{itemize}

\sources
\begin{itemize}
    \item IETF, RFC 4786 : Operation of Anycast Services,
    \url{https://www.rfc-editor.org/rfc/rfc4786}.
\end{itemize}

\subsubsection*{Question 57 (Pratique) : Haute disponibilité DNS avec plusieurs datacenters}

\reponse
\begin{enumerate}
    \item \textbf{DNS autoritaire redondant :} au moins deux
    fournisseurs DNS ou un service Anycast, avec des serveurs de noms
    dans plusieurs régions.
    \item \textbf{GeoDNS :} chaque région est dirigée vers son
    datacenter le plus proche (Europe \fl{} DC Paris, Afrique \fl{} DC
    Alger).
    \item \textbf{Health checks} (HTTP, TCP) sur chaque datacenter, avec
    \textbf{DNS Failover} automatique vers un autre site si l'un tombe.
    \item \textbf{TTL court} (60 s) sur les enregistrements concernés,
    pour une bascule rapide.
    \item \textbf{Pondération} pour répartir la charge ou faire une
    migration progressive.
    \item \textbf{DNSSEC} pour protéger les réponses.
    \item Dans chaque datacenter, des load balancers locaux redondants,
    et des données répliquées entre les sites.
\end{enumerate}

\exemple
Avec AWS Route 53 : politique de routage par géolocalisation, combinée
à une politique de basculement et à des health checks sur chaque
région.

\sources
\begin{itemize}
    \item AWS, Amazon Route 53 Developer Guide,
    \url{https://docs.aws.amazon.com/route53/}.
\end{itemize}

% --------------------------------------------------
\subsection{Membre 9 : API REST, microservices et WebSocket}
% --------------------------------------------------
\membre{\dotfill}{58 à 62}

\paragraph{F. API REST et Microservices}

\subsubsection*{Question 58 (Facile) : Qu'est-ce qu'une API REST ?}

\reponse
Une API REST est une interface qui permet à des applications de
communiquer via HTTP, en manipulant des \textbf{ressources} identifiées
par des URI (\texttt{/etudiants/12}), le plus souvent en JSON.

\textbf{Principes fondamentaux (Roy Fielding, 2000) :}
\begin{enumerate}
    \item \textbf{Client-serveur :} séparation entre l'interface et le
    stockage des données.
    \item \textbf{Sans état (stateless) :} chaque requête contient tout
    ce qu'il faut pour être traitée, le serveur ne garde pas de session.
    \item \textbf{Mise en cache :} les réponses indiquent si elles
    peuvent être mises en cache.
    \item \textbf{Interface uniforme :} URI, méthodes HTTP standards
    (GET, POST, PUT, DELETE), codes de statut.
    \item \textbf{Système en couches :} des intermédiaires (proxy, load
    balancer) peuvent s'insérer de façon transparente.
    \item \textbf{Code à la demande} (optionnel) : le serveur peut
    envoyer du code exécutable.
\end{enumerate}

\sources
\begin{itemize}
    \item R. Fielding, Architectural Styles and the Design of
    Network-based Software Architectures, thèse, University of
    California, Irvine, 2000.
\end{itemize}

\subsubsection*{Question 59 (Moyen) : Services réseau essentiels dans une architecture microservices}

\reponse
\begin{itemize}
    \item \textbf{Découverte de services :} un registre (Consul, Eureka,
    etcd) ou le DNS interne de Kubernetes (CoreDNS) permet à chaque
    service de trouver les autres, même quand leurs IP changent.
    \item \textbf{Communication inter-services :} synchrone (REST, gRPC)
    ou asynchrone (files de messages : Kafka, RabbitMQ). Un
    \textbf{service mesh} (Istio, Linkerd) peut gérer le chiffrement
    mTLS, le routage et la répartition de charge.
    \item \textbf{Monitoring et tracing distribué :} métriques
    (Prometheus, Grafana), logs centralisés (ELK), traçage d'une requête
    à travers tous les services (OpenTelemetry, Jaeger).
    \item \textbf{Gestion des échecs :}
    \begin{itemize}
        \item \textbf{Retry} avec attente exponentielle
        (\textit{exponential backoff}) pour les erreurs temporaires ;
        \item \textbf{Circuit Breaker} : après plusieurs échecs, les
        appels vers le service défaillant sont coupés pendant un temps,
        pour éviter les pannes en cascade (Resilience4j, Istio) ;
        \item \textbf{timeouts} et réponses de secours (\textit{fallback}).
    \end{itemize}
\end{itemize}

\sources
\begin{itemize}
    \item M. Fowler, CircuitBreaker,
    \url{https://martinfowler.com/bliki/CircuitBreaker.html}.
\end{itemize}

\subsubsection*{Question 60 (Avancé) : Rôle d'une API Gateway}

\reponse
L'API Gateway est le \textbf{point d'entrée unique} des clients vers
les microservices. Elle cache la complexité interne : le client
n'appelle qu'une seule adresse.

\textbf{Fonctionnalités attendues :}
\begin{itemize}
    \item routage des requêtes vers le bon microservice ;
    \item authentification et autorisation (OAuth2, JWT, clés d'API) ;
    \item limitation de débit (\textit{rate limiting}) et quotas ;
    \item terminaison TLS ;
    \item agrégation de plusieurs appels en une seule réponse ;
    \item transformation de protocoles et de formats (REST vers gRPC) ;
    \item cache, journalisation, métriques et traçage ;
    \item gestion des versions d'API.
\end{itemize}

\exemple
Kong, NGINX, Spring Cloud Gateway, AWS API Gateway.

\sources
\begin{itemize}
    \item C. Richardson, Pattern: API Gateway,
    \url{https://microservices.io/patterns/apigateway.html}.
\end{itemize}

\paragraph{G. WebSocket et Communication en Temps Réel}

\subsubsection*{Question 61 (Moyen) : HTTP classique et WebSocket}

\reponse
\begin{itemize}
    \item \textbf{HTTP classique :} modèle requête/réponse. Seul le
    client peut initier un échange, et chaque requête porte des
    en-têtes complets. Le serveur ne peut pas envoyer de données de
    lui-même.
    \item \textbf{WebSocket :} après une négociation HTTP
    (\texttt{Upgrade: websocket}), la connexion TCP reste
    \textbf{ouverte}, \textbf{bidirectionnelle} et \textbf{full-duplex}.
    Client et serveur envoient des messages à tout moment, avec très
    peu de surcharge (ws:// ou wss:// avec TLS).
\end{itemize}

\textbf{Cas où WebSocket est préférable :} chat et messagerie
instantanée, jeux en ligne, tableaux de bord temps réel, trading,
édition collaborative, notifications instantanées.

\sources
\begin{itemize}
    \item IETF, RFC 6455 : The WebSocket Protocol,
    \url{https://www.rfc-editor.org/rfc/rfc6455}.
\end{itemize}

\subsubsection*{Question 62 (Pratique) : Comparaison des technologies temps réel}

\reponse
\begin{longtable}{|L{2.6cm}|L{6.2cm}|L{5.8cm}|}
\hline
\textbf{Technologie} & \textbf{Fonctionnement} & \textbf{Cas d'usage
approprié} \\ \hline
\endhead
WebSocket & Connexion persistante, bidirectionnelle. & Chat, jeux
multijoueurs, édition collaborative. \\ \hline
Server-Sent Events (SSE) & Flux HTTP persistant, du serveur vers le
client uniquement, reconnexion automatique. & Fil d'actualités, cours
de bourse, notifications, suivi de commande. \\ \hline
Long Polling & Le client envoie une requête que le serveur garde ouverte
jusqu'à avoir une donnée, puis le client relance. & Solution de repli
pour les anciens navigateurs ou les réseaux qui bloquent WebSocket. \\
\hline
HTTP/2 Server Push & Le serveur envoie des ressources liées à une
requête avant qu'elles soient demandées. & Précharger CSS et JS d'une
page. Ce n'est pas un vrai canal temps réel, et il est aujourd'hui
abandonné par les navigateurs. \\ \hline
\end{longtable}

\sources
\begin{itemize}
    \item WHATWG, HTML Living Standard : Server-sent events,
    \url{https://html.spec.whatwg.org/multipage/server-sent-events.html}.
    \item IETF, RFC 6455, \url{https://www.rfc-editor.org/rfc/rfc6455}.
\end{itemize}

% --------------------------------------------------
\subsection{Membre 10 : Routage avancé}
% --------------------------------------------------
\membre{Abdelmadjid Baghdali \& Ouweis Wail Sari Bey}{63 à 71}

\paragraph{A. BGP (Border Gateway Protocol)}

\subsubsection*{Question 63 (Facile) : BGP, protocole de routage inter-domaines}

\reponse
BGP est \textbf{inter-domaines} car il route le trafic entre des
\textbf{AS (Systèmes Autonomes)} gérés par des organisations
indépendantes. Son rôle principal est d'\textbf{échanger les préfixes
IP joignables entre AS}, ce qui fait fonctionner \textbf{Internet à
l'échelle mondiale}.

\subsubsection*{Question 64 (Moyen) : Différence entre iBGP et eBGP}

\reponse
\begin{itemize}
    \item \textbf{eBGP :} session BGP entre routeurs de \textbf{deux AS
    différents}, généralement directement connectés. Elle sert à
    échanger les routes entre AS, et l'ASN est ajouté à
    l'\textbf{AS\_PATH}.
    \item \textbf{iBGP :} session BGP entre routeurs du \textbf{même AS},
    pas forcément voisins. Elle sert à distribuer à l'intérieur de
    l'AS les routes apprises en eBGP. L'\textbf{AS\_PATH} et le
    \textbf{NEXT\_HOP} ne sont pas modifiés, et il faut un
    \textbf{full mesh} ou un \textbf{Route Reflector}.
    \item \textbf{Différence principale :} eBGP fonctionne
    \textbf{entre AS} (échange externe), alors qu'iBGP fonctionne
    \textbf{dans un même AS} (distribution interne).
\end{itemize}

\subsubsection*{Question 65 (Avancé) : Les Systèmes Autonomes (AS)}

\reponse
\begin{itemize}
    \item \textbf{C'est quoi un AS ?}
    Un groupe de réseaux/routeurs sous une même autorité
    administrative, avec une seule politique de routage vers Internet.
    \item \textbf{Comment les AS sont-ils identifiés ?}
    Par un \textbf{ASN} (numéro de 16 ou 32 bits), attribué par
    l'\textbf{IANA} via les \textbf{RIR}.
    \item \textbf{Pourquoi sont-ils importants pour BGP ?}
    BGP route \textbf{entre AS}, pas entre routeurs individuels.
    L'\textbf{AS\_PATH} permet d'éviter les boucles et de choisir les
    chemins, et les ASN servent à appliquer les politiques de routage
    et à vérifier l'origine des préfixes (\textbf{RPKI}).
\end{itemize}

\subsubsection*{Question 66 (Réflexion) : Scénario : entreprise multi-homée (2 FAI)}

\reponse
Une entreprise e-commerce avec son propre \textbf{ASN} et un bloc
\textbf{PI} se connecte à \textbf{deux FAI} via \textbf{eBGP}.

\textbf{Pourquoi BGP est indispensable :}
\begin{itemize}
    \item \textbf{Redondance :} si le lien FAI 1 tombe, BGP
    \textbf{retire la route} et bascule vers le FAI 2, sans changer d'IP.
    \item \textbf{Contrôle du trafic :} \textbf{local preference}
    (sortant), \textbf{AS-path prepending} / \textbf{communautés}
    (entrant).
    \item \textbf{Indépendance :} adresses \textbf{PI}, donc pas de
    renumérotation en changeant de FAI.
    \item \textbf{Sécurité :} \textbf{RPKI/ROA} contre les
    \textbf{BGP hijacks}.
\end{itemize}

\paragraph{B. OSPF (Open Shortest Path First)}

\subsubsection*{Question 67 (Moyen) : Avantages d'OSPF par rapport à RIP}

\reponse
\begin{itemize}
    \item \textbf{Type :} OSPF est \textbf{à état de liens} (algorithme
    de \textbf{Dijkstra / SPF}), RIP est \textbf{à vecteur de distance}.
    \item \textbf{Métrique :} OSPF utilise le \textbf{coût} (basé sur la
    bande passante), RIP le \textbf{nombre de sauts} (max 15).
    \item \textbf{Convergence :} OSPF converge \textbf{rapidement} grâce
    aux mises à jour déclenchées ; RIP envoie toute sa table toutes les
    30 s et converge lentement.
    \item \textbf{Évolutivité :} OSPF gère les grands réseaux grâce aux
    \textbf{areas} ; RIP est limité à 15 sauts.
    \item \textbf{Autres :} VLSM/CIDR, authentification, pas de boucles
    de routage, moins de trafic de mise à jour.
\end{itemize}

\subsubsection*{Question 68 (Avancé) : Les areas OSPF et l'Area 0}

\reponse
\begin{itemize}
    \item \textbf{Area :} groupe de routeurs qui partagent la même base
    de données d'état de liens (\textbf{LSDB}). Découper en areas
    \textbf{limite l'inondation des LSA}, réduit la taille des tables et
    le calcul SPF, et permet de \textbf{résumer} les routes.
    \item \textbf{Area 0 (backbone) :} zone centrale. \textbf{Toutes les
    autres areas doivent y être connectées} (sinon : \textbf{virtual
    link}). Tout le trafic inter-areas passe par l'Area 0.
\end{itemize}

\subsubsection*{Question 69 (Pratique) : Types de routeurs OSPF}

\reponse
\begin{itemize}
    \item \textbf{DR (Designated Router) :} élu sur un réseau
    multi-accès (Ethernet), il centralise les échanges d'état de liens,
    ce qui évite un full mesh d'adjacences. Élection : \textbf{priorité}
    la plus haute, puis \textbf{Router-ID} le plus haut.
    \item \textbf{BDR (Backup DR) :} remplace le DR en cas de panne.
    \item \textbf{ABR (Area Border Router) :} relie \textbf{une area à
    l'Area 0} et résume les routes entre areas.
    \item \textbf{ASBR (AS Boundary Router) :} \textbf{injecte des routes
    externes} dans OSPF (statiques, BGP, autre protocole).
\end{itemize}

\paragraph{C. MPLS (Multiprotocol Label Switching)}

\subsubsection*{Question 70 (Moyen) : Principe de fonctionnement de MPLS}

\reponse
MPLS achemine les paquets selon un \textbf{label} court (20 bits) placé
entre l'en-tête de couche 2 et l'en-tête IP (« couche 2,5 »), au lieu
de chercher l'adresse IP à chaque saut.
\begin{itemize}
    \item Le \textbf{LER d'entrée} classe le paquet dans une
    \textbf{FEC} et lui \textbf{ajoute un label} (push).
    \item Les \textbf{LSR} du cœur \textbf{remplacent le label} (swap)
    en suivant un chemin prédéfini, le \textbf{LSP}.
    \item Le \textbf{LER de sortie} \textbf{retire le label} (pop) et
    livre le paquet IP.
    \item Les labels sont distribués par \textbf{LDP} ou
    \textbf{RSVP-TE}.
\end{itemize}

\subsubsection*{Question 71 (Avancé) : Avantages de MPLS}

\reponse
\begin{itemize}
    \item \textbf{VPN :} les \textbf{L3VPN} (VRF + MP-BGP, RFC 4364)
    isolent le trafic de plusieurs clients sur la même infrastructure,
    même avec des adresses privées identiques. MPLS permet aussi des
    \textbf{L2VPN} (VPLS, pseudowires).
    \item \textbf{Ingénierie de trafic :} \textbf{RSVP-TE} impose des
    chemins explicites, indépendants du plus court chemin IGP, pour
    \textbf{équilibrer les liens} et réserver de la bande passante.
    \textbf{Fast Reroute} bascule en moins de 50 ms.
    \item \textbf{Qualité de service :} les \textbf{bits TC/EXP} (3 bits)
    du label portent la classe de trafic, compatible
    \textbf{DiffServ} ; la voix et la vidéo sont prioritaires de bout en
    bout.
\end{itemize}

\sources
\begin{itemize}
    \item IETF, RFC 4271 (BGP-4), RFC 2328 (OSPFv2), RFC 3031 (MPLS
    Architecture), RFC 4364 (BGP/MPLS IP VPNs),
    \url{https://www.rfc-editor.org/}.
\end{itemize}

% --------------------------------------------------
\subsection{Membre 11 : IPsec et SSL/TLS}
% --------------------------------------------------
\membre{\dotfill}{72 à 77}

\paragraph{A. IPsec}

\subsubsection*{Question 72 (Facile) : Modes transport et tunnel d'IPsec}

\reponse
\begin{itemize}
    \item \textbf{IPsec mode transport:} protects the useful part of
    the IP packs, and eventually some parts of the header but it
    conserves the original header of the pack, it's generally used in
    communication host to host.
    \item \textbf{IPsec mode tunnel:} it encapsulates the entire IP pack
    with its header included (header + payload), in a new IP pack with
    a new header so the original pack will be the payload for the new
    pack with a new outer IP header, its used for VPN site to site or
    gateway to gateway, because it allows the transition of the trafic
    between two distant networks in a secured way.
\end{itemize}

\subsubsection*{Question 73 (Moyen) : AH et ESP}

\reponse
\begin{itemize}
    \item \textbf{AH} provides only authentication and integrity for
    packets, with no encryption or confidentiality. It protects the IP
    packet including the IP header, except for mutable fields, so the
    outer IP header is authenticated. However, this causes problems
    with NAT, because NAT modifies the IP header that AH authenticates.
    AH provides authentication without confidentiality and corresponds
    to IP protocol number 51.
    \item \textbf{ESP} can provide authentication, integrity, and
    confidentiality. It protects the payload and, in tunnel mode, the
    inner IP packet, but the outer IP header is not authenticated by
    ESP. ESP is more NAT-friendly than AH, and NAT traversal is handled
    through NAT-T. It is mostly used in IPsec VPNs to protect data
    through encryption and corresponds to IP protocol number 50.
\end{itemize}

\subsubsection*{Question 74 (Pratique) : VPN IPsec site-à-site}

\reponse
\textbf{Scénario :} une entreprise a un siège à Alger et une agence à
Oran, et veut relier leurs réseaux internes à travers Internet, de
façon sécurisée et sans louer de ligne dédiée.

\textbf{Architecture type :} chaque site possède une passerelle VPN
(routeur ou pare-feu) avec une IP publique. Un tunnel IPsec en
\textbf{mode tunnel} avec \textbf{ESP} est établi entre les deux
passerelles. Les postes n'ont rien à configurer : le trafic entre
192.168.10.0/24 (Alger) et 192.168.20.0/24 (Oran) est chiffré
automatiquement par les passerelles.

L'établissement se fait avec \textbf{IKEv2} (UDP 500, et 4500 avec
NAT-T) : authentification par clé pré-partagée ou certificats, échange
Diffie-Hellman, puis création des SA IPsec. Paramètres conseillés :
AES-256-GCM, SHA-256, groupe DH 19 ou plus.

\paragraph{B. SSL/TLS}

\subsubsection*{Question 75 (Facile) : Différence entre SSL et TLS}

\reponse
SSL a été créé par Netscape (SSL 2.0 en 1995, SSL 3.0 en 1996). TLS est
son successeur standardisé par l'IETF : TLS 1.0 (1999), 1.1 (2006), 1.2
(2008) et 1.3 (2018). TLS corrige les failles de SSL (POODLE, etc.) et
utilise des algorithmes plus robustes.

SSL 1.0, 2.0, 3.0, TLS 1.0 et TLS 1.1 sont aujourd'hui \textbf{obsolètes
et vulnérables}. TLS 1.2 reste acceptable, mais \textbf{TLS 1.3 est
recommandé} : plus sûr, plus rapide, et la \textit{forward secrecy} y est
obligatoire.

Le terme « SSL » est encore souvent utilisé par abus de langage
(« certificat SSL »), mais ce sont des connexions \textbf{TLS} qui sont
déployées aujourd'hui.

\subsubsection*{Question 76 (Moyen) : Le handshake TLS}

\reponse
Le handshake permet de \textbf{négocier les paramètres
cryptographiques}, d'\textbf{authentifier le serveur} et de
\textbf{dériver les clés de session}.

\textbf{TLS 1.2 (environ 2 allers-retours) :} ClientHello \fl{}
ServerHello \fl{} Certificate \fl{} échange de clés (ECDHE) \fl{}
Finished.

\textbf{TLS 1.3 (1-RTT) :} messages superflus supprimés, forward secrecy
obligatoire.
\begin{enumerate}
    \item \textbf{ClientHello :} le client envoie les versions et suites
    de chiffrement supportées, un nombre aléatoire et sa part de clé
    Diffie-Hellman (\textit{key\_share}).
    \item \textbf{ServerHello :} le serveur choisit la suite et envoie
    sa part de clé. Les deux calculent alors la même clé secrète.
    \item Le serveur envoie, déjà chiffrés, son \textbf{certificat}, une
    signature \textbf{CertificateVerify} (preuve qu'il possède la clé
    privée) et un message \textbf{Finished}.
    \item Le client vérifie le certificat, puis envoie son
    \textbf{Finished}.
    \item Les données applicatives circulent chiffrées avec des clés
    symétriques (AES-GCM, ChaCha20).
\end{enumerate}

\subsubsection*{Question 77 (Réflexion) : Certificats numériques et autorités de certification}

\reponse
Les certificats sont essentiels pour \textbf{lier une clé publique à une
identité} (nom de domaine) et empêcher les attaques de type
\textbf{homme du milieu}. Sans certificat valide, le chiffrement peut
exister, mais \textbf{sans authentification fiable} : le client ne sait
pas s'il parle au vrai serveur ou à un attaquant. Le \textbf{certificat
X.509} contient la clé publique, l'identité et une date de validité.

L'\textbf{autorité de certification (CA)} est un tiers de confiance qui
vérifie l'identité du demandeur, \textbf{signe} le certificat et gère
la \textbf{révocation} (CRL, OCSP). Le navigateur fait confiance aux CA
racines déjà installées et vérifie la \textbf{chaîne de confiance}
(racine \fl{} intermédiaire \fl{} certificat du site). La CA permet ainsi
à un client de faire confiance à un serveur qu'il ne connaît pas a
priori. Exemples : Let's Encrypt, DigiCert.

\textbf{En bref :} TLS 1.3 est la version à utiliser ; le handshake
établit une session sécurisée en 1-RTT ; les certificats et la CA
garantissent l'authenticité du serveur et fondent la confiance sur
Internet.

\sources
\begin{itemize}
    \item IETF, RFC 4301 (IPsec), RFC 7296 (IKEv2), RFC 8446 (TLS 1.3),
    RFC 8996 (dépréciation de TLS 1.0/1.1), RFC 5280 (certificats X.509),
    \url{https://www.rfc-editor.org/}.
\end{itemize}

% --------------------------------------------------
\subsection{Membre 12 : Firewalls et NAT}
% --------------------------------------------------
\membre{Benbraham Sara}{78 et 79}

\paragraph{C. Firewalls et NAT}

\subsubsection*{Question 78 (Moyen) : Différences entre les types de firewalls}

\reponse
Un firewall (pare-feu) est un dispositif matériel ou logiciel qui
contrôle le trafic réseau selon des règles de sécurité afin
d'autoriser ou de bloquer les communications.

\begin{enumerate}
    \item \textbf{Stateless (filtrage de paquets)}

    Chaque paquet est examiné indépendamment, principalement selon les
    adresses IP, les ports, le protocole et éventuellement les flags
    TCP, à l'aide de règles de filtrage (ACL). Il est rapide et peu
    coûteux en ressources, mais ne conserve aucun état des connexions.
    Une règle trop permissive peut donc laisser passer des paquets
    forgés qui semblent appartenir à un trafic autorisé.

    \item \textbf{Stateful (inspection avec état)}

    Le firewall maintient une table d'états contenant les informations
    sur les connexions actives (adresses IP, ports, protocole, état et
    temporisation). Pour une nouvelle connexion TCP, il vérifie
    généralement le paquet initial SYN selon ses règles, puis autorise
    les paquets de retour correspondant à la connexion suivie. Il offre
    un meilleur contrôle contextuel, mais consomme davantage de
    ressources et peut être ciblé par une attaque SYN flood visant à
    épuiser ses ressources.

    \item \textbf{Firewall applicatif (proxy, couche 7 du modèle OSI)}

    Le proxy termine la connexion du client et établit une seconde
    connexion vers le serveur. Il peut analyser les requêtes
    applicatives (URL, commandes et contenu), appliquer des règles
    d'authentification et bloquer certaines requêtes malveillantes. Il
    offre un contrôle plus approfondi, mais peut introduire de la
    latence et nécessite davantage de ressources. L'inspection du
    contenu HTTPS nécessite généralement un mécanisme de déchiffrement
    approprié.
\end{enumerate}

\textbf{Exemples :}
\begin{itemize}
    \item \textbf{Stateless :} autoriser les paquets TCP à destination
    du port 80.
    \item \textbf{Stateful :} autoriser les réponses appartenant à une
    connexion TCP préalablement établie et suivie.
    \item \textbf{Proxy :} bloquer une requête HTTP vers une URL
    interdite ou détectée comme malveillante.
\end{itemize}

\subsubsection*{Question 79 (Avancé) : Fonctionnement du NAT et ses limitations}

\reponse
\textbf{1. Définition}

Le NAT (Network Address Translation) est un mécanisme qui modifie les
adresses IP dans les paquets réseau lors de leur passage entre deux
réseaux. Il permet notamment à plusieurs appareils d'un réseau privé
d'accéder à Internet en partageant une adresse IPv4 publique, lorsque
le NAT est associé à la traduction des ports (PAT).

\textbf{2. Fonctionnement du NAT}

Le processus se déroule en quatre étapes :
\begin{enumerate}
    \item L'ordinateur envoie un paquet contenant son adresse IP privée
    comme adresse source.
    \item Le routeur remplace l'adresse source privée par son adresse
    IP publique. Avec le PAT, il traduit également le port source.
    \item Le serveur distant envoie sa réponse à l'adresse IP publique
    et au port traduits.
    \item Le routeur consulte sa table de traduction pour retrouver
    l'adresse IP et le port privés du destinataire, puis lui transmet
    la réponse.
\end{enumerate}

Le \textbf{PAT} (Port Address Translation), également appelé
\textbf{NAT overload}, permet à plusieurs appareils de partager une
même adresse IPv4 publique en utilisant des ports distincts pour
différencier leurs communications.

\textbf{3. Les limitations du NAT}
\begin{itemize}
    \item \textbf{Connexions entrantes :} un appareil derrière un NAT
    n'est généralement pas directement accessible depuis Internet sans
    configuration spécifique, comme une redirection de ports.
    \item \textbf{Compatibilité des protocoles :} certains protocoles
    transportent des adresses IP ou des ports dans leurs données et
    nécessitent des mécanismes supplémentaires pour fonctionner
    correctement à travers le NAT.
    \item \textbf{Communications de bout en bout :} le NAT complique
    les connexions directes entre appareils, notamment pour certaines
    applications P2P et VoIP.
    \item \textbf{Épuisement des ports :} avec le PAT, les ports
    disponibles et les ressources de traduction peuvent devenir des
    contraintes en cas de nombreuses connexions simultanées.
    \item \textbf{Dépannage plus difficile :} les traductions
    d'adresses compliquent l'identification des appareils internes et
    l'analyse des journaux réseau.
    \item \textbf{Sécurité limitée :} le NAT n'est pas un mécanisme de
    sécurité suffisant à lui seul. Il ne remplace pas un firewall
    configuré avec des règles de filtrage.
\end{itemize}

\sources
\begin{itemize}
    \item RFC 2979 ; RFC 4987 ; RFC 3022 ; RFC 4787 (voir la
    bibliographie générale).
\end{itemize}

% --------------------------------------------------
\subsection{Membre 13 : QoS}
% --------------------------------------------------
\membre{\dotfill}{80 à 83}

\subsubsection*{Question 80 (Facile) : Pourquoi la QoS est-elle nécessaire ?}

\reponse
Un réseau transporte en même temps des trafics aux besoins très
différents. Sans QoS, tous les paquets sont traités de la même façon
(\textit{best effort}) : en cas de congestion, un téléchargement peut
dégrader un appel VoIP ou une visioconférence. La QoS garantit aux
flux critiques une \textbf{bande passante}, une \textbf{latence}, une
\textbf{gigue} et un \textbf{taux de perte} adaptés.

\exemple
La voix exige une latence de moins de 150 ms (ITU-T G.114), une gigue
inférieure à 30 ms et moins de 1~\% de perte.

\sources
\begin{itemize}
    \item ITU-T, Recommandation G.114 : One-way transmission time,
    \url{https://www.itu.int/}.
\end{itemize}

\subsubsection*{Question 81 (Moyen) : Classification, marquage, mise en file d'attente, policing et shaping}

\reponse
\begin{itemize}
    \item \textbf{Classification :} identifier le type de trafic (ACL,
    ports, adresses, NBAR/DPI).
    \item \textbf{Marquage :} inscrire la classe dans le paquet
    (\textbf{DSCP} dans l'en-tête IP, \textbf{CoS} 802.1p en couche 2),
    pour que les équipements suivants n'aient pas à reclasser.
    \item \textbf{Mise en file d'attente :} en cas de congestion, les
    paquets sont placés dans des files selon leur classe :
    \textbf{LLQ} (file prioritaire stricte pour la voix), \textbf{CBWFQ}
    (bande passante garantie par classe), \textbf{WRED} (rejet anticipé
    pour éviter la saturation).
    \item \textbf{Policing :} limite le débit en \textbf{rejetant} (ou
    re-marquant) l'excédent, sans mise en mémoire.
    \item \textbf{Shaping :} limite le débit en \textbf{retardant}
    l'excédent dans un tampon ; le trafic est lissé, au prix d'un peu de
    latence.
\end{itemize}

\sources
\begin{itemize}
    \item IETF, RFC 2475 : An Architecture for Differentiated Services,
    \url{https://www.rfc-editor.org/rfc/rfc2475}.
\end{itemize}

\subsubsection*{Question 82 (Pratique) : Priorisation de la voix, de la vidéo et des données}

\reponse
\begin{longtable}{|L{2.4cm}|L{2.4cm}|L{3.6cm}|L{6.2cm}|}
\hline
\textbf{Priorité} & \textbf{Trafic} & \textbf{Marquage} &
\textbf{Justification} \\ \hline
1 & Voix & EF (DSCP 46), file LLQ & Très sensible à la latence et à la
gigue, débit faible et constant. Limitée à environ 33~\% du lien pour
ne pas affamer les autres. \\ \hline
2 & Vidéo & AF41 (DSCP 34), CBWFQ & Sensible à la latence, débit élevé ;
bande passante garantie. \\ \hline
3 & Données critiques & AF21 / AF31 & Applications métier, tolèrent un
peu de délai grâce aux retransmissions TCP. \\ \hline
4 & Autres données & Best effort (DSCP 0) & Navigation, e-mail,
téléchargements : élastiques et peu sensibles. \\ \hline
\end{longtable}

\sources
\begin{itemize}
    \item IETF, RFC 4594 : Configuration Guidelines for DiffServ Service
    Classes, \url{https://www.rfc-editor.org/rfc/rfc4594}.
\end{itemize}

\subsubsection*{Question 83 (Avancé) : DiffServ et IntServ}

\reponse
\begin{longtable}{|L{3.2cm}|L{5.7cm}|L{5.7cm}|}
\hline
\textbf{Critère} & \textbf{IntServ} & \textbf{DiffServ} \\ \hline
\endhead
Principe & Réservation de ressources \textbf{par flux} & Traitement
\textbf{par classe} (PHB) \\ \hline
Signalisation & RSVP, de bout en bout & Aucune : marquage DSCP \\ \hline
État dans les routeurs & Un état par flux & Aucun état par flux \\ \hline
Garantie & Stricte & Relative (statistique) \\ \hline
Évolutivité & Faible (impossible sur Internet) & Très bonne \\ \hline
\end{longtable}

\textbf{DiffServ est le plus utilisé aujourd'hui}, car il passe à
l'échelle : les routeurs du cœur n'ont qu'un petit nombre de classes à
gérer, quel que soit le nombre de flux. IntServ reste limité à de
petits réseaux ou combiné à MPLS-TE.

\sources
\begin{itemize}
    \item IETF, RFC 1633 (IntServ), RFC 2205 (RSVP), RFC 2474 et
    RFC 2475 (DiffServ), \url{https://www.rfc-editor.org/}.
\end{itemize}

% --------------------------------------------------
\subsection{Membre 14 : SDN et NFV}
% --------------------------------------------------
\membre{\dotfill}{84 à 88}

\paragraph{A. SDN (Software-Defined Networking)}

\subsubsection*{Question 84 : Qu'est-ce que le SDN ? Séparation du plan de contrôle et du plan de données}

\reponse
Le SDN (Software-Defined Networking) est une architecture réseau qui
sépare l'intelligence du réseau (les décisions d'acheminement) des
équipements qui transportent réellement les paquets. Cette intelligence
est centralisée dans un contrôleur logiciel programmable, qui dispose
d'une vue globale de la topologie.

\begin{itemize}
    \item \textbf{Plan de données (data plane) :} il transfère les
    paquets (commutation, routage, filtrage) en appliquant des règles
    déjà définies. Il est exécuté par les switchs et routeurs, à la
    vitesse du matériel.
    \item \textbf{Plan de contrôle (control plane) :} il décide où
    envoyer le trafic (calcul des routes, tables MAC, politiques). Dans
    un réseau traditionnel, chaque équipement possède son propre plan de
    contrôle (OSPF, STP, etc.) ; en SDN, il est centralisé dans le
    contrôleur (OpenDaylight, ONOS, Ryu, etc.).
\end{itemize}

Une couche applications (pare-feu, répartition de charge, QoS,
supervision) s'ajoute au-dessus du contrôleur. L'architecture comporte
donc trois couches reliées par deux types d'interfaces :

\begin{longtable}{|L{3cm}|L{5.8cm}|L{5.8cm}|}
\hline
\textbf{Couche} & \textbf{Rôle} & \textbf{Interface} \\ \hline
Applications & Pare-feu, load balancer, QoS, supervision & API
northbound (REST) vers le contrôleur \\ \hline
Contrôleur & Cerveau logique, vue globale du réseau & API southbound
(OpenFlow, NETCONF) vers les équipements \\ \hline
Infrastructure & Switchs et routeurs qui transfèrent les paquets & -- \\
\hline
\end{longtable}

\exemple
Dans un réseau d'entreprise, un administrateur définit depuis le
contrôleur une politique unique (« le VLAN invités ne peut pas accéder
aux serveurs internes »). Le contrôleur la traduit en règles installées
automatiquement sur tous les switchs concernés, sans configurer chaque
équipement un par un.

\sources
\begin{itemize}
    \item IETF, RFC 7426 : Software-Defined Networking (SDN) : Layers and
    Architecture Terminology, \url{https://www.rfc-editor.org/rfc/rfc7426},
    consulté le 18 octobre 2026.
    \item Open Networking Foundation, SDN Architecture Overview,
    \url{https://opennetworking.org/}, consulté le 18 octobre 2026.
\end{itemize}

\subsubsection*{Question 85 : Qu'est-ce qu'OpenFlow ? Quel est son rôle dans le SDN ?}

\reponse
OpenFlow est un protocole standardisé par l'ONF (Open Networking
Foundation). C'est le protocole southbound historique du SDN : il permet
au contrôleur de programmer directement les tables de flux des
équipements réseau.

\begin{itemize}
    \item Un switch OpenFlow contient une ou plusieurs \textbf{tables de
    flux}. Chaque entrée comprend des champs de correspondance (match :
    port d'entrée, adresses MAC et IP, VLAN, ports TCP/UDP, etc.), des
    actions (transférer, rejeter, modifier un en-tête, envoyer au
    contrôleur), des compteurs et des délais d'expiration.
    \item Le canal entre le contrôleur et le switch est sécurisé (TLS).
    \item Messages principaux : \textbf{Packet-In} (le switch ne connaît
    pas le flux et interroge le contrôleur), \textbf{Flow-Mod} (le
    contrôleur installe ou modifie une règle), \textbf{Packet-Out} et les
    messages de statistiques.
\end{itemize}

\textbf{Rôle dans le SDN :} OpenFlow est le « langage commun » entre le
plan de contrôle centralisé et le plan de données. Il permet un contrôle
fin, flux par flux, indépendant du constructeur des équipements.

\exemple
Un premier paquet d'un nouveau flux arrive sur un switch qui n'a aucune
règle correspondante. Le switch envoie un Packet-In au contrôleur ;
celui-ci décide du chemin et répond par un Flow-Mod. Les paquets
suivants du même flux sont ensuite traités directement par le switch,
sans solliciter à nouveau le contrôleur.

\sources
\begin{itemize}
    \item Open Networking Foundation, OpenFlow Switch Specification,
    \url{https://opennetworking.org/}, consulté le 18 octobre 2026.
    \item IETF, RFC 7426, \url{https://www.rfc-editor.org/rfc/rfc7426},
    consulté le 18 octobre 2026.
\end{itemize}

\subsubsection*{Question 86 : Quels avantages le SDN apporte-t-il par rapport aux réseaux traditionnels ?}

\reponse
\begin{itemize}
    \item \textbf{Centralisation et vue globale :} une politique
    cohérente est appliquée à l'ensemble du réseau depuis un point
    unique.
    \item \textbf{Programmabilité et automatisation :} la configuration
    se fait par API ou scripts, ce qui évite la configuration manuelle
    équipement par équipement et réduit les erreurs humaines.
    \item \textbf{Agilité :} le déploiement de VLAN, de services ou de
    politiques passe de plusieurs jours à quelques minutes.
    \item \textbf{Ouverture :} les standards ouverts réduisent la
    dépendance à un constructeur unique.
    \item \textbf{Optimisation du trafic :} l'ingénierie de trafic
    dynamique permet une meilleure utilisation de la bande passante.
    \item \textbf{Sécurité :} isolation et micro-segmentation
    dynamiques, mise en quarantaine automatique d'un hôte compromis.
    \item \textbf{Réduction des coûts} d'investissement et
    d'exploitation.
\end{itemize}

\textbf{Limites à signaler :} le contrôleur est un point de défaillance
unique s'il n'est pas déployé en cluster, il devient une cible
d'attaque privilégiée, et la migration depuis un réseau existant est
complexe.

\exemple
\begin{longtable}{|L{3.6cm}|L{5.5cm}|L{5.5cm}|}
\hline
\textbf{Critère} & \textbf{Réseau traditionnel} & \textbf{Réseau SDN}
\\ \hline
Configuration & Manuelle, équipement par équipement & Centralisée,
automatisée \\ \hline
Plan de contrôle & Distribué dans chaque équipement & Centralisé dans le
contrôleur \\ \hline
Évolutivité des politiques & Lente, risque d'incohérences & Rapide,
cohérente \\ \hline
Dépendance constructeur & Forte & Réduite (standards ouverts) \\ \hline
\end{longtable}

\sources
\begin{itemize}
    \item Cisco, Software-Defined Networking (SDN) : vue d'ensemble,
    \url{https://www.cisco.com/}, consulté le 18 octobre 2026.
    \item Open Networking Foundation, \url{https://opennetworking.org/},
    consulté le 18 octobre 2026.
\end{itemize}

\paragraph{B. NFV (Network Function Virtualization)}

\subsubsection*{Question 87 : Qu'est-ce que la NFV ? Exemples de fonctions réseau virtualisables}

\reponse
La NFV (Network Function Virtualization) consiste à remplacer les
équipements réseau dédiés (boîtiers matériels propriétaires) par des
fonctions réseau logicielles, appelées \textbf{VNF} (Virtualized Network
Functions), qui s'exécutent sur des serveurs standards via des machines
virtuelles ou des conteneurs.

L'architecture définie par l'ETSI comporte trois blocs :
\begin{itemize}
    \item \textbf{VNF :} la fonction réseau logicielle elle-même ;
    \item \textbf{NFVI} (NFV Infrastructure) : serveurs, hyperviseur et
    réseau virtuel qui hébergent les VNF ;
    \item \textbf{MANO} (Management and Orchestration) : gestion du cycle
    de vie et orchestration des VNF.
\end{itemize}

\textbf{Exemples de fonctions virtualisables :}
\begin{itemize}
    \item routeur virtuel (vRouter) et switch virtuel (Open vSwitch) ;
    \item pare-feu et NGFW virtuels ;
    \item IDS/IPS (Snort, Suricata) ;
    \item répartiteur de charge, reverse proxy et WAF ;
    \item NAT, DHCP et DNS ;
    \item passerelle VPN et SD-WAN ;
    \item DPI et optimiseur WAN ;
    \item cœur de réseau mobile (vEPC, cœur 5G) et IMS.
\end{itemize}

\exemple
Au lieu d'acheter un pare-feu matériel pour chaque site, une entreprise
déploie un pare-feu virtuel sur un serveur standard. En cas de pic de
trafic, l'orchestrateur lance une seconde instance en quelques minutes.

\sources
\begin{itemize}
    \item ETSI, GS NFV 002 : Network Functions Virtualisation (NFV) ;
    Architectural Framework, \url{https://www.etsi.org/technologies/nfv},
    consulté le 18 octobre 2026.
\end{itemize}

\subsubsection*{Question 88 : Relation entre SDN et NFV : complémentaires ou concurrents ?}

\reponse
Le SDN et la NFV sont \textbf{complémentaires et non concurrents}. Ils
sont indépendants (on peut déployer l'un sans l'autre), mais leur
association est particulièrement efficace.

\begin{longtable}{|L{2.4cm}|L{6.1cm}|L{6.1cm}|}
\hline
 & \textbf{SDN} & \textbf{NFV} \\ \hline
Objectif & Contrôler et programmer le réseau (séparer contrôle et
données) & Virtualiser les fonctions réseau (découpler logiciel et
matériel) \\ \hline
Cible & L'acheminement du trafic & Les fonctions (pare-feu, load
balancer, etc.) \\ \hline
Moyen & Contrôleur et protocole OpenFlow & VNF sur serveurs standards \\
\hline
\end{longtable}

\textbf{Synergie :} la NFV fournit des fonctions souples et élastiques ;
le SDN dirige dynamiquement le trafic à travers elles. On parle de
\textbf{service chaining} : par exemple client \fl{} pare-feu virtuel
\fl{} IDS virtuel \fl{} répartiteur de charge virtuel. L'orchestrateur
crée ou supprime les VNF selon la charge, et le contrôleur SDN adapte
les chemins en conséquence.

\exemple
Dans un datacenter, un contrôleur SDN oriente le trafic web entrant vers
un WAF virtuel, puis vers un load balancer virtuel, avant d'atteindre
les serveurs d'application. Chaque VNF peut être dupliquée ou retirée
sans modifier la topologie physique.

\sources
\begin{itemize}
    \item ETSI, GS NFV 002, \url{https://www.etsi.org/technologies/nfv},
    consulté le 18 octobre 2026.
    \item IETF, RFC 7426, \url{https://www.rfc-editor.org/rfc/rfc7426},
    consulté le 18 octobre 2026.
\end{itemize}

% --------------------------------------------------
\subsection{Membre 15 : Synthèse globale et architecture}
% --------------------------------------------------
\membre{\dotfill}{89}

\subsubsection*{Question 89 : Infrastructure réseau d'un campus de 5000 utilisateurs}

\textbf{Énoncé :} concevoir l'infrastructure réseau d'un campus
universitaire de 5000 utilisateurs, avec 3 bâtiments distants, des
services administratifs, pédagogiques et de recherche, des besoins de
sécurité élevés et une connexion Internet via deux FAI différents.

\paragraph{a) Services réseau obligatoires et justification}

\begin{longtable}{|L{4.4cm}|L{10.2cm}|}
\hline
\textbf{Service} & \textbf{Justification} \\ \hline
\endhead
VLAN et trunking 802.1Q & Segmenter les populations (administratif,
pédagogique, recherche, invités, serveurs, gestion), limiter les
domaines de broadcast et isoler les zones. \\ \hline
Routage inter-VLAN (switchs L3 au cœur) & Routage interne performant et
application des ACL entre zones. \\ \hline
OSPF (interne) et BGP (vers les FAI) & Convergence rapide, évolutivité
et gestion du multi-homing avec deux FAI. \\ \hline
DHCP (avec failover) & Adressage automatique de 5000 utilisateurs, avec
un pool par VLAN. \\ \hline
DNS (interne, externe, secondaire) & Résolution de noms, indispensable à
tous les services ; split-DNS entre interne et externe. \\ \hline
NAT/PAT & Économie d'adresses publiques et masquage du plan d'adressage
interne. \\ \hline
Pare-feu NGFW et DMZ & Filtrage périmétrique et protection des serveurs
exposés. \\ \hline
IDS/IPS & Détection et blocage des attaques. \\ \hline
AAA (RADIUS/802.1X, TACACS+) & Authentification, autorisation et
traçabilité des accès des utilisateurs et des administrateurs. \\ \hline
Annuaire (Active Directory/LDAP) & Gestion centralisée des identités. \\
\hline
VPN (IPsec/SSL) & Accès distant sécurisé pour les enseignants et
chercheurs, interconnexion de sites. \\ \hline
NTP & Horloges synchronisées (journaux, certificats, Kerberos). \\ \hline
Syslog/SIEM, SNMP, NetFlow & Supervision, corrélation d'événements et
audit. \\ \hline
Proxy, reverse proxy, WAF & Filtrage web sortant, protection et
publication des applications. \\ \hline
QoS & Priorité à la voix, à la visioconférence et aux applications
critiques. \\ \hline
Sauvegarde des configurations & Restauration rapide après un incident
(SFTP/TFTP sécurisé). \\ \hline
WiFi (contrôleurs WLC, WPA2/WPA3-Enterprise) & Mobilité des 5000
utilisateurs. \\ \hline
\end{longtable}

\paragraph{b) Authentification et gestion des utilisateurs}

\begin{itemize}
    \item \textbf{Annuaire central AD/LDAP} répliqué (au moins deux
    contrôleurs de domaine par site principal), avec des groupes :
    étudiants, enseignants, administratifs, chercheurs, invités.
    \item \textbf{802.1X avec RADIUS} (FreeRADIUS ou NPS) sur les ports
    filaires et le WiFi, en EAP-TLS ou PEAP, couplé à l'annuaire. Le
    serveur RADIUS attribue dynamiquement le VLAN selon le groupe de
    l'utilisateur. Une fédération de type eduroam permet l'accueil des
    visiteurs académiques.
    \item \textbf{MFA} pour les comptes sensibles (administrateurs,
    recherche, VPN).
    \item \textbf{TACACS+} pour l'accès aux équipements réseau, avec
    contrôle d'accès basé sur les rôles (RBAC) et journalisation de
    chaque commande.
    \item \textbf{Portail captif} pour les invités : comptes temporaires,
    VLAN isolé, accès à Internet uniquement.
    \item \textbf{Cycle de vie des comptes :} création automatique à
    l'inscription, désactivation au départ, politique de mots de passe,
    audit périodique.
\end{itemize}

\paragraph{c) Supervision et monitoring}

\begin{itemize}
    \item \textbf{Disponibilité et performance :} Zabbix, LibreNMS ou
    PRTG via SNMPv3 (CPU, mémoire, interfaces, température), avec seuils
    et alertes par mail ou SMS.
    \item \textbf{Analyse du trafic :} NetFlow/IPFIX (ntopng, Elastiflow)
    pour détecter les anomalies et planifier la capacité.
    \item \textbf{Journalisation centralisée :} Syslog vers une pile
    ELK/SIEM (Elasticsearch, Logstash, Kibana) avec des règles de
    détection (force brute, scans, mouvements latéraux), idéalement
    opérée par un SOC.
    \item \textbf{Tableaux de bord} Grafana ou Kibana, sondes actives
    (ping, HTTP, DNS) et mesure des SLA.
    \item \textbf{Réseau de gestion séparé} (VLAN dédié ou
    \textit{out-of-band}) pour que la supervision reste disponible en
    cas de panne du réseau de production.
\end{itemize}

\paragraph{d) Redondance et haute disponibilité}

\begin{itemize}
    \item \textbf{Topologie hiérarchique} (cœur, distribution, accès)
    avec un double cœur ; chaque bâtiment est relié aux deux cœurs.
    \item \textbf{Liaisons inter-bâtiments en fibre redondante}, par des
    chemins physiques distincts.
    \item \textbf{EtherChannel (LACP)} pour agréger les liens et tolérer
    une panne, \textbf{RSTP/MST} contre les boucles.
    \item \textbf{FHRP (HSRP/VRRP)} : passerelle par défaut redondante
    pour chaque VLAN.
    \item \textbf{Double FAI :} BGP (idéalement avec un bloc d'adresses
    et un AS propres) ou, à défaut, routes statiques flottantes avec IP
    SLA, voire routage par politique ; fonctionnement actif/passif ou
    avec répartition de charge.
    \item \textbf{Pare-feu en cluster HA} (actif/passif avec
    synchronisation d'état), alimentations redondantes et onduleurs.
    \item \textbf{Services critiques répliqués :} DHCP failover, DNS
    maître et esclaves, plusieurs contrôleurs AD et serveurs RADIUS,
    virtualisation en cluster.
    \item \textbf{Sauvegardes testées} et plan de reprise d'activité
    (PRA).
\end{itemize}

\paragraph{e) Mécanismes de sécurité par niveau}

\begin{longtable}{|L{3.4cm}|L{11.2cm}|}
\hline
\textbf{Niveau} & \textbf{Mesures} \\ \hline
\endhead
Physique & Locaux techniques fermés, contrôle d'accès,
vidéosurveillance, ports inutilisés désactivés. \\ \hline
Couche 2 (accès) & 802.1X, port security, DHCP snooping, DAI, IP Source
Guard, BPDU guard et root guard, VLAN natif inutilisé, désactivation de
DTP. \\ \hline
Couche 3 (réseau) & Segmentation par VLAN et zones, ACL entre zones
(recherche isolée, invités sans accès interne), uRPF, authentification
OSPF/BGP. \\ \hline
Périmètre & NGFW, IPS, DMZ, protection anti-DDoS, filtrage URL et DNS.
\\ \hline
Application & WAF, reverse proxy, TLS partout, durcissement des
serveurs, mises à jour régulières. \\ \hline
Terminaux & Antivirus/EDR, contrôle de conformité (NAC), chiffrement
des disques. \\ \hline
Administration & SSH uniquement, TACACS+, réseau de gestion isolé,
sauvegarde des configurations. \\ \hline
Détection et réponse & SIEM, alertes, procédures de gestion d'incident,
audits et tests d'intrusion. \\ \hline
Données & VPN IPsec/SSL, chiffrement, politique de sauvegarde. \\ \hline
\end{longtable}

\textbf{Principes directeurs :} défense en profondeur, moindre
privilège et Zero Trust.

\paragraph{f) Gestion de la QoS}

\begin{enumerate}
    \item \textbf{Classification et marquage} à la frontière de
    confiance (switchs d'accès) en DSCP ; les téléphones IP et
    équipements de confiance conservent leur marquage.
    \item \textbf{Classes de trafic :}
    \begin{itemize}
        \item voix (VoIP) : EF, file prioritaire stricte (LLQ), faible
        latence ;
        \item visioconférence et cours en ligne : AF4x, bande passante
        garantie ;
        \item applications critiques (administration, ENT,
        authentification) : AF3x ;
        \item trafic standard : best effort ;
        \item trafic non essentiel (P2P, streaming, téléchargements
        massifs) : scavenger (CS1), bande passante limitée.
    \end{itemize}
    \item \textbf{Mécanismes :} files CBWFQ/LLQ, WRED contre la
    congestion, policing en entrée et shaping sur les liens vers les FAI.
    \item \textbf{Cohérence :} même politique sur tous les équipements
    (cœur, distribution, WAN), avec un suivi de la bande passante par FAI
    via NetFlow.
\end{enumerate}

\paragraph{g) Schéma simplifié de l'architecture}

\begin{center}
\begin{tikzpicture}[
  box/.style={draw, rounded corners, align=center, font=\small,
              minimum width=3.2cm, minimum height=1cm, fill=blue!5},
  >=latex]
  \node[box] (f1) at (-2.5,6) {FAI 1\\BGP, lien 1};
  \node[box] (f2) at (2.5,6)  {FAI 2\\BGP, lien 2};
  \node[box, fill=red!8] (fw) at (0,4.3) {Pare-feu NGFW (HA)\\IPS, NAT, VPN};
  \node[box] (dmz) at (5,4.3) {DMZ\\Reverse proxy, WAF};
  \node[box] (c1) at (-2.2,2.5) {Cœur 1 (L3)\\OSPF, HSRP};
  \node[box] (c2) at (2.2,2.5)  {Cœur 2 (L3)\\OSPF, HSRP};
  \node[box] (sup) at (-6,2.5) {Supervision\\SIEM, Zabbix};
  \node[box] (dc) at (6,2.5) {Datacenter\\AD, DNS, DHCP};
  \node[box] (ba) at (-4.2,0.3) {Bâtiment A\\Accès + WiFi};
  \node[box] (bb) at (0,0.3)    {Bâtiment B\\Accès + WiFi};
  \node[box] (bc) at (4.2,0.3)  {Bâtiment C\\Accès + WiFi};
  \draw (f1)--(fw); \draw (f2)--(fw); \draw (fw)--(dmz);
  \draw (fw)--(c1); \draw (fw)--(c2);
  \draw[very thick] (c1)--node[above,font=\scriptsize]{LACP}(c2);
  \draw (sup)--(c1); \draw (c2)--(dc);
  \foreach \b in {ba,bb,bc}{\draw (\b)--(c1); \draw (\b)--(c2);}
\end{tikzpicture}
\end{center}

Chaque bâtiment est relié aux deux cœurs par des liaisons en fibre
redondantes (LACP). Les VLAN (administratif, pédagogique, recherche,
invités, serveurs, gestion) sont définis dans les trois bâtiments et
routés par les cœurs, avec des ACL entre zones. La QoS est appliquée de
bout en bout.

\paragraph{h) Intérêt des technologies web et d'infrastructure}

\textbf{1. Proxy et reverse proxy.}

\begin{longtable}{|L{2.6cm}|L{6cm}|L{6cm}|}
\hline
 & \textbf{Proxy (forward)} & \textbf{Reverse proxy} \\ \hline
Position & Devant les clients (côté utilisateur) & Devant les serveurs
(côté application) \\ \hline
Rôle & Les clients internes accèdent à Internet par son intermédiaire &
Les clients externes accèdent aux serveurs par son intermédiaire \\ \hline
Cas d'usage & Filtrage web, contrôle d'accès, journalisation,
anonymisation & Publication d'applications, terminaison TLS, répartition
de charge \\ \hline
Exemples & Squid & Nginx, HAProxy, Traefik \\ \hline
\end{longtable}

\begin{itemize}
    \item \textbf{Sécurité :} le proxy filtre les URL, bloque les
    logiciels malveillants et trace les accès ; le reverse proxy masque
    l'architecture interne, centralise TLS, intègre un WAF et limite les
    attaques (DDoS, injections).
    \item \textbf{Performance :} cache des contenus, compression,
    mutualisation des connexions, terminaison TLS déchargée des serveurs.
    \item \textbf{Disponibilité :} le reverse proxy répartit la charge,
    détecte les serveurs en panne (health checks) et bascule le trafic.
\end{itemize}

\textbf{2. CDN (Content Delivery Network).}
\begin{itemize}
    \item \textbf{Principe :} réseau de serveurs de cache (edge servers)
    répartis géographiquement. L'utilisateur est dirigé, via DNS ou
    Anycast, vers le serveur le plus proche, qui sert le contenu sans
    solliciter le serveur d'origine. Cela réduit la latence et la charge
    sur l'origine, et absorbe les pics de trafic et une partie des
    attaques DDoS.
    \item \textbf{Stratégies de cache :}
    \begin{itemize}
        \item durée de vie (TTL) et en-têtes Cache-Control et Expires ;
        \item revalidation conditionnelle avec ETag et Last-Modified
        (réponse 304) ;
        \item contenu statique (images, CSS, JS) mis en cache longtemps
        avec des noms de fichiers versionnés ; contenu dynamique peu ou
        pas mis en cache ;
        \item purge ou invalidation manuelle ou automatique ; clé de
        cache (URL, en-têtes, cookies) ;
        \item modèles pull (le CDN récupère à la demande) ou push
        (contenu préchargé).
    \end{itemize}
\end{itemize}
Exemples : Cloudflare, Akamai, Amazon CloudFront.

\textbf{3. Load balancing web.}
\begin{itemize}
    \item \textbf{Intérêt :} répartir les requêtes sur plusieurs
    serveurs pour la performance, la scalabilité et la disponibilité
    (suppression des points de défaillance uniques), et permettre la
    maintenance sans interruption.
    \item \textbf{Niveaux :} couche 4 (TCP/UDP, rapide) et couche 7
    (HTTP : routage selon l'URL, les en-têtes ou les cookies).
    \item \textbf{Algorithmes :} round robin, round robin pondéré, least
    connections, IP hash, least response time.
    \item \textbf{Fonctions :} health checks, persistance de session,
    terminaison SSL, redondance du répartiteur lui-même
    (Keepalived/VRRP).
\end{itemize}
Exemples : HAProxy, Nginx, F5, AWS Elastic Load Balancing.

\textbf{4. DNS pour services web.}
\begin{itemize}
    \item \textbf{GeoDNS :} répond avec une adresse IP différente selon
    la localisation du client, afin de diriger l'utilisateur vers le
    datacenter le plus proche (latence, conformité).
    \item \textbf{DNS Failover :} des sondes surveillent la santé des
    serveurs et le DNS retire automatiquement l'adresse défaillante, avec
    un TTL court pour une bascule rapide ; limite : dépend des caches
    DNS.
    \item \textbf{Anycast :} la même adresse IP est annoncée en BGP
    depuis plusieurs sites ; le routage achemine le client vers le site
    le plus proche. Il offre une faible latence, de la résilience et une
    bonne résistance aux DDoS (serveurs DNS racine, CDN).
\end{itemize}

\textbf{5. API REST et microservices.}
\begin{itemize}
    \item \textbf{Principes REST :} architecture client-serveur, sans
    état (stateless), interface uniforme (ressources identifiées par des
    URI, verbes HTTP GET, POST, PUT, DELETE, représentations en JSON),
    réponses pouvant être mises en cache, système en couches, codes de
    statut HTTP standards.
    \item \textbf{Architecture microservices :} l'application est
    découpée en petits services indépendants, chacun avec une fonction
    métier et sa propre base de données, déployables et dimensionnables
    séparément (conteneurs, Docker, Kubernetes). Avantages : agilité,
    résilience, scalabilité ciblée. Inconvénients : complexité, latence
    réseau, cohérence des données, supervision plus difficile.
    \item \textbf{API Gateway :} point d'entrée unique devant les
    microservices. Elle assure le routage, l'authentification et
    l'autorisation (OAuth2, JWT), la limitation de débit (rate limiting),
    le cache, la composition de réponses, la journalisation et la
    supervision (Kong, NGINX, AWS API Gateway).
\end{itemize}

\textbf{6. WebSocket.} Communication temps réel : après une négociation
HTTP (\texttt{Upgrade: websocket}), WebSocket établit une connexion TCP
persistante et bidirectionnelle (\textit{full-duplex}) entre le client et
le serveur. Chaque partie peut envoyer des messages à tout moment, avec
très peu d'en-têtes. Usages : chat, jeux en ligne, trading, collaboration
en direct, notifications.

\begin{longtable}{|L{2.2cm}|L{4.1cm}|L{4.1cm}|L{4.1cm}|}
\hline
 & \textbf{WebSocket} & \textbf{SSE (Server-Sent Events)} &
\textbf{Long Polling} \\ \hline
Sens & Bidirectionnel & Serveur vers client uniquement & Pseudo-push
(requêtes répétées) \\ \hline
Connexion & Persistante (ws/wss) & Persistante (HTTP) & Requête
maintenue puis renouvelée \\ \hline
Complexité & Plus élevée & Simple, reconnexion automatique & Simple,
compatible partout \\ \hline
Surcharge & Faible & Faible & Élevée (en-têtes HTTP à chaque cycle) \\
\hline
Cas d'usage & Chat, jeux, temps réel interactif & Flux de notifications,
cours de bourse & Solution de repli, anciens navigateurs \\ \hline
\end{longtable}

\sources
\begin{itemize}
    \item IETF, RFC 2474 : Definition of the Differentiated Services
    Field, \url{https://www.rfc-editor.org/rfc/rfc2474}, consulté le 18
    octobre 2026.
    \item IETF, RFC 2865 : RADIUS,
    \url{https://www.rfc-editor.org/rfc/rfc2865}, consulté le 18 octobre
    2026.
    \item IETF, RFC 4786 : Operation of Anycast Services,
    \url{https://www.rfc-editor.org/rfc/rfc4786}, consulté le 18 octobre
    2026.
    \item IETF, RFC 6455 : The WebSocket Protocol,
    \url{https://www.rfc-editor.org/rfc/rfc6455}, consulté le 18 octobre
    2026.
    \item IETF, RFC 9110 : HTTP Semantics,
    \url{https://www.rfc-editor.org/rfc/rfc9110}, consulté le 18 octobre
    2026.
    \item Cisco, documentation réseau, \url{https://www.cisco.com/},
    consulté le 18 octobre 2026.
\end{itemize}

% ==================================================
\section{Conclusion générale}
% ==================================================

Les services étudiés forment un ensemble cohérent. Les services
fondamentaux (DNS, DHCP, NTP) assurent la connectivité de base ;
l'annuaire et AAA (LDAP, RADIUS, TACACS+, Kerberos) contrôlent qui accède
à quoi ; la supervision (SNMP, Syslog, NetFlow) donne la visibilité ; le
routage (OSPF, BGP, MPLS) relie les sites et Internet ; la sécurité
(IPsec, TLS, pare-feu, NAT) protège les échanges ; la QoS garantit la
performance des flux critiques ; enfin, les services web, le SDN et la
NFV apportent la disponibilité, l'automatisation et la flexibilité
nécessaires à une infrastructure d'entreprise moderne.

% ==================================================
\section{Bibliographie générale}
% ==================================================

\begin{enumerate}
    \item RFC 1034, RFC 1035 (DNS) ; RFC 4033 (DNSSEC).
    \item RFC 2131, RFC 2132 (DHCP) ; RFC 5905 (NTPv4).
    \item RFC 4511, RFC 4512, RFC 4514 (LDAP) ; RFC 2865, RFC 2866
    (RADIUS) ; RFC 8907 (TACACS+) ; RFC 4120 (Kerberos V5).
    \item RFC 1213, RFC 3411, RFC 3414, RFC 3416 (SNMP) ; RFC 5424
    (Syslog) ; RFC 3954 (NetFlow v9) ; RFC 7011 (IPFIX).
    \item RFC 9110, RFC 9111, RFC 9113, RFC 9114 (HTTP) ; RFC 9000 (QUIC)
    ; RFC 5789 (PATCH) ; RFC 1794 (DNS load balancing) ; RFC 4786
    (Anycast) ; RFC 6455 (WebSocket).
    \item R. Fielding, Architectural Styles and the Design of
    Network-based Software Architectures, thèse, 2000.
    \item RFC 4271 (BGP-4) ; RFC 1930 (AS) ; RFC 6793 (ASN 4 octets) ;
    RFC 2328 (OSPFv2) ; RFC 3031 (MPLS) ; RFC 4364 (BGP/MPLS IP VPNs).
    \item RFC 4301, RFC 4302, RFC 4303 (IPsec) ; RFC 7296 (IKEv2) ;
    RFC 8446 (TLS 1.3) ; RFC 8996 ; RFC 5280 (X.509).
    \item RFC 2979 -- Behavior of and Requirements for Internet
    Firewalls ; RFC 4987 -- TCP SYN Flooding Attacks and Common
    Mitigations.
    \item RFC 3022 -- Traditional IP Network Address Translator ;
    RFC 4787 -- Network Address Translation Behavioral Requirements for
    UDP.
    \item RFC 1633 (IntServ) ; RFC 2205 (RSVP) ; RFC 2474, RFC 2475,
    RFC 4594 (DiffServ) ; ITU-T G.114.
    \item RFC 7426 (SDN) ; ONF, OpenFlow Switch Specification ; ETSI GS
    NFV 002.
\end{enumerate}

\end{document}